\documentclass[12pt,a4paper]{article}

\usepackage{amsmath,amssymb,amsfonts}
\usepackage{bm}
\usepackage{physics}
\usepackage{graphicx}
\usepackage{hyperref}
\usepackage{geometry}
\geometry{margin=1in}

\title{Special Relativistic Magnetohydrodynamics in Godunov SPH Format:\\
A Complete Mathematical Formulation from First Principles}

\author{Based on Kitajima et al. (2025) and Iwasaki \& Inutsuka (2011)}

\date{\today}

\begin{document}

\maketitle

\begin{abstract}
This document presents a comprehensive mathematical formulation for Special Relativistic
Magnetohydrodynamics (SRMHD) using the Godunov Smoothed Particle Hydrodynamics (GSPH) framework.
We derive all equations from first principles, beginning with special relativity and Maxwell's
equations, proceeding through the stress-energy tensor construction, and culminating in the
full GSPH discretization. Particular attention is given to the primitive variable recovery
problem, where the presence of magnetic fields creates a coupled nonlinear system requiring
numerical iteration. We combine the special relativistic Godunov SPH method of Kitajima et al. (2025)
with the MHD extension of Iwasaki \& Inutsuka (2011) to derive a complete set of equations
suitable for numerical implementation.
\end{abstract}

\tableofcontents
\newpage

%=============================================================================
\section{Introduction}
%=============================================================================

Special Relativistic Magnetohydrodynamics (SRMHD) is essential for simulating high-energy
astrophysical phenomena such as relativistic jets, gamma-ray bursts, and compact object
mergers. This document presents a Godunov SPH formulation that:

\begin{enumerate}
    \item Uses Riemann solvers to capture shocks without artificial viscosity
    \item Handles relativistic effects through proper Lorentz transformations
    \item Incorporates magnetic fields via the MHD stress tensor
    \item Uses the Method of Characteristics for Alfv\'en waves
    \item Maintains conservation laws through symmetric particle interactions
\end{enumerate}

\subsection{Notation and Conventions}

Throughout this document, we use the following conventions:
\begin{itemize}
    \item Speed of light: $c$ (often set to 1 in natural units)
    \item Metric signature: $(-,+,+,+)$ (Minkowski spacetime)
    \item Greek indices $\alpha, \beta, \mu, \nu$: spacetime components (0, 1, 2, 3)
    \item Latin indices $i, j, k$: spatial components (1, 2, 3)
    \item Bold symbols: 3-vectors in space, e.g., $\bm{v}$, $\bm{B}$
    \item Subscripts $L$, $R$: left and right states in Riemann problem
    \item Superscript $*$: interface (star region) values from Riemann solution
    \item Subscript $\parallel$, $\perp$: parallel and perpendicular to direction of interest
    \item Einstein summation convention: repeated indices are summed
\end{itemize}

%=============================================================================
\section{Foundations of Special Relativity}
%=============================================================================

\subsection{Spacetime and Four-Vectors}

In special relativity, spacetime is described by Minkowski geometry with metric:
\begin{equation}
\eta_{\mu\nu} = \text{diag}(-1, +1, +1, +1)
\label{eq:minkowski_metric}
\end{equation}

An event is specified by the four-vector $x^\mu = (ct, x, y, z) = (ct, \bm{x})$.

\subsubsection{Four-Velocity}

For a particle or fluid element moving with three-velocity $\bm{v}$, the four-velocity is:
\begin{equation}
u^\mu = \gamma(c, \bm{v}) = \gamma(c, v^x, v^y, v^z)
\label{eq:four_velocity}
\end{equation}

where the Lorentz factor $\gamma$ is:
\begin{equation}
\gamma = \frac{1}{\sqrt{1 - v^2/c^2}} = \frac{1}{\sqrt{1 - \beta^2}}, \quad \beta = \frac{|\bm{v}|}{c}
\label{eq:lorentz_factor}
\end{equation}

The four-velocity satisfies the normalization:
\begin{equation}
u^\mu u_\mu = \eta_{\mu\nu} u^\mu u^\nu = -\gamma^2 c^2 + \gamma^2 v^2 = -c^2
\label{eq:four_velocity_norm}
\end{equation}

This normalization is a fundamental constraint that will be important for variable recovery.

\subsubsection{Proper Time and Derivatives}

The proper time $\tau$ is the time measured in the rest frame of the fluid element:
\begin{equation}
d\tau = \frac{dt}{\gamma}
\label{eq:proper_time}
\end{equation}

The comoving (Lagrangian) derivative is:
\begin{equation}
\frac{D}{D\tau} = u^\mu \partial_\mu = \gamma\left(\frac{\partial}{\partial t} + \bm{v} \cdot \nabla\right)
\label{eq:comoving_derivative}
\end{equation}

\subsection{Relativistic Thermodynamics}

\subsubsection{Rest-Frame Quantities}

In the fluid rest frame, we define:
\begin{itemize}
    \item $\rho_0$: rest-mass density (mass per unit proper volume)
    \item $n$: baryon number density (baryons per unit proper volume)
    \item $\epsilon$: specific internal energy (energy per unit rest mass)
    \item $p$: thermodynamic pressure
    \item $T$: temperature
    \item $s$: specific entropy
\end{itemize}

The rest-mass density and baryon density are related by:
\begin{equation}
\rho_0 = m_B n
\label{eq:rest_mass_baryon}
\end{equation}
where $m_B$ is the mean baryon mass.

\subsubsection{Relativistic Enthalpy}

The specific enthalpy $h$ (enthalpy per unit rest mass) is:
\begin{equation}
h = c^2 + \epsilon + \frac{p}{\rho_0} = c^2\left(1 + \frac{\epsilon}{c^2} + \frac{p}{\rho_0 c^2}\right)
\label{eq:specific_enthalpy}
\end{equation}

In units where $c = 1$:
\begin{equation}
h = 1 + \epsilon + \frac{p}{\rho_0}
\label{eq:specific_enthalpy_c1}
\end{equation}

The total enthalpy density (enthalpy per unit proper volume) is:
\begin{equation}
w = \rho_0 h = \rho_0 c^2 + \rho_0 \epsilon + p = e + p
\label{eq:enthalpy_density}
\end{equation}
where $e = \rho_0 c^2 + \rho_0 \epsilon$ is the total energy density in the rest frame.

\subsubsection{Equation of State}

For an ideal gas with adiabatic index $\Gamma$:
\begin{equation}
p = (\Gamma - 1)\rho_0 \epsilon
\label{eq:ideal_gas_eos}
\end{equation}

The specific enthalpy becomes:
\begin{equation}
h = c^2 + \epsilon + \frac{(\Gamma-1)\epsilon}{1} = c^2 + \frac{\Gamma}{\Gamma-1}\frac{p}{\rho_0}
\label{eq:enthalpy_ideal_gas}
\end{equation}

In units where $c = 1$:
\begin{equation}
h = 1 + \frac{\Gamma}{\Gamma-1}\frac{p}{\rho_0}
\label{eq:enthalpy_ideal_gas_c1}
\end{equation}

The relativistic sound speed is derived from:
\begin{equation}
c_s^2 = \left.\frac{\partial p}{\partial e}\right|_s = \frac{\Gamma p}{w} = \frac{\Gamma p}{\rho_0 h}
\label{eq:sound_speed_derivation}
\end{equation}

For the ideal gas EOS:
\begin{equation}
c_s^2 = \frac{\Gamma(\Gamma-1)(h-c^2)}{\Gamma h - (\Gamma-1)c^2}
\label{eq:sound_speed_ideal}
\end{equation}

In the non-relativistic limit $h \to c^2$, this gives $c_s^2 \to \Gamma p/\rho_0$.

%=============================================================================
\section{Electromagnetic Fields in Special Relativity}
%=============================================================================

\subsection{Electromagnetic Field Tensor}

The electric field $\bm{E}$ and magnetic field $\bm{B}$ are components of the
antisymmetric electromagnetic field tensor $F^{\mu\nu}$:
\begin{equation}
F^{\mu\nu} = \begin{pmatrix}
0 & -E^x/c & -E^y/c & -E^z/c \\
E^x/c & 0 & -B^z & B^y \\
E^y/c & B^z & 0 & -B^x \\
E^z/c & -B^y & B^x & 0
\end{pmatrix}
\label{eq:em_tensor}
\end{equation}

The dual tensor is:
\begin{equation}
{}^*F^{\mu\nu} = \frac{1}{2}\epsilon^{\mu\nu\alpha\beta}F_{\alpha\beta} = \begin{pmatrix}
0 & -B^x & -B^y & -B^z \\
B^x & 0 & E^z/c & -E^y/c \\
B^y & -E^z/c & 0 & E^x/c \\
B^z & E^y/c & -E^x/c & 0
\end{pmatrix}
\label{eq:dual_em_tensor}
\end{equation}

\subsection{Maxwell's Equations in Covariant Form}

Maxwell's equations in vacuum become:
\begin{align}
\partial_\nu F^{\mu\nu} &= \mu_0 J^\mu \label{eq:maxwell_1} \\
\partial_\nu {}^*F^{\mu\nu} &= 0 \label{eq:maxwell_2}
\end{align}

Equation \eqref{eq:maxwell_2} gives $\nabla \cdot \bm{B} = 0$ and Faraday's law.
Equation \eqref{eq:maxwell_1} gives Gauss's law and Amp\`ere's law.

\subsection{Magnetic Field Four-Vector}

In ideal MHD, we work in the fluid rest frame where the electric field vanishes
(perfect conductivity). The magnetic field four-vector $b^\mu$ is defined such that:
\begin{equation}
b^\mu u_\mu = 0
\label{eq:b_orthogonal}
\end{equation}

The components are:
\begin{align}
b^0 &= \frac{\gamma}{c}(\bm{v} \cdot \bm{B}) = \frac{W}{c}v^i B_i \label{eq:b0} \\
b^i &= \frac{B^i}{\gamma} + \frac{\gamma}{c^2}(\bm{v} \cdot \bm{B})v^i = \frac{B^i + (b^0/c)v^i}{\gamma}
\label{eq:bi}
\end{align}

where $W = \gamma$ is the Lorentz factor (notation used in some references).

The magnitude squared of the magnetic four-vector:
\begin{equation}
b^2 = b^\mu b_\mu = \frac{B^2}{\gamma^2} + \frac{(\bm{v} \cdot \bm{B})^2}{c^2}
= \frac{B^2 + (\bm{v} \cdot \bm{B})^2/c^2}{\gamma^2}
\label{eq:b_squared}
\end{equation}

In the fluid rest frame, $b^2 = B'^2$ where $B'$ is the magnetic field measured
in the comoving frame.

\subsection{Ideal MHD Condition: Derivation from Perfect Conductivity}

The ideal MHD condition is one of the most important assumptions in MHD. We derive
it from first principles by considering the limit of infinite electrical conductivity.

\subsubsection{Ohm's Law in a Conductor}

In a stationary conductor, Ohm's law relates the current density $\bm{J}$ to the
electric field $\bm{E}$:
\begin{equation}
\bm{J} = \sigma \bm{E}
\label{eq:ohms_law_stationary}
\end{equation}
where $\sigma$ is the electrical conductivity (units: S/m or $\Omega^{-1}$m$^{-1}$).

\subsubsection{Ohm's Law in a Moving Conductor}

For a conductor moving with velocity $\bm{v}$ through a magnetic field $\bm{B}$,
charges experience both the electric field and the Lorentz force. The generalized
Ohm's law becomes:
\begin{equation}
\bm{J} = \sigma (\bm{E} + \bm{v} \times \bm{B})
\label{eq:ohms_law_moving}
\end{equation}

The term $\bm{E} + \bm{v} \times \bm{B}$ is precisely the electric field as measured
in the rest frame of the moving fluid (to first order in $v/c$). This can be seen
from the Lorentz transformation of electromagnetic fields.

\subsubsection{Lorentz Transformation of Electric Field: Full Derivation}

We derive how electromagnetic fields transform under Lorentz boosts using the
field tensor formalism.

\textbf{Step 1: The Lorentz Transformation Matrix}

Consider a boost along an arbitrary direction $\hat{\bm{n}} = \bm{v}/|\bm{v}|$
with velocity $\bm{v}$. The Lorentz transformation matrix $\Lambda^\mu_{\ \nu}$
transforms coordinates as $x'^\mu = \Lambda^\mu_{\ \nu} x^\nu$.

For a boost along $\hat{\bm{n}}$:
\begin{equation}
\Lambda^\mu_{\ \nu} = \begin{pmatrix}
\gamma & -\gamma \beta n_x & -\gamma \beta n_y & -\gamma \beta n_z \\
-\gamma \beta n_x & 1 + (\gamma-1)n_x^2 & (\gamma-1)n_x n_y & (\gamma-1)n_x n_z \\
-\gamma \beta n_y & (\gamma-1)n_y n_x & 1 + (\gamma-1)n_y^2 & (\gamma-1)n_y n_z \\
-\gamma \beta n_z & (\gamma-1)n_z n_x & (\gamma-1)n_z n_y & 1 + (\gamma-1)n_z^2
\end{pmatrix}
\label{eq:lorentz_matrix}
\end{equation}

where $\beta = v/c$ and $\gamma = (1-\beta^2)^{-1/2}$.

The spatial part can be written compactly as:
\begin{equation}
\Lambda^i_{\ j} = \delta^i_j + (\gamma - 1)\frac{v^i v_j}{v^2}
\label{eq:lorentz_spatial}
\end{equation}

\textbf{Step 2: Transformation of the Field Tensor}

The electromagnetic field tensor $F^{\mu\nu}$ transforms as a rank-2 tensor:
\begin{equation}
F'^{\mu\nu} = \Lambda^\mu_{\ \alpha} \Lambda^\nu_{\ \beta} F^{\alpha\beta}
\label{eq:tensor_transform}
\end{equation}

\textbf{Step 3: Extract Electric Field Components}

The electric field components are $E^i = c F^{0i}$, so:
\begin{equation}
E'^i = c F'^{0i} = c \Lambda^0_{\ \alpha} \Lambda^i_{\ \beta} F^{\alpha\beta}
\label{eq:Ei_transform_start}
\end{equation}

Expanding the sum over $\alpha$ and $\beta$:
\begin{align}
E'^i &= c \left[ \Lambda^0_{\ 0} \Lambda^i_{\ j} F^{0j}
+ \Lambda^0_{\ k} \Lambda^i_{\ 0} F^{k0}
+ \Lambda^0_{\ 0} \Lambda^i_{\ 0} F^{00}
+ \Lambda^0_{\ k} \Lambda^i_{\ l} F^{kl} \right]
\end{align}

Using $F^{00} = 0$, $F^{0j} = E^j/c$, $F^{k0} = -E^k/c$, and
$F^{kl} = -\epsilon^{klm} B_m$:
\begin{align}
E'^i &= \Lambda^0_{\ 0} \Lambda^i_{\ j} E^j
- \Lambda^0_{\ k} \Lambda^i_{\ 0} E^k
- c \Lambda^0_{\ k} \Lambda^i_{\ l} \epsilon^{klm} B_m
\label{eq:Ei_expanded}
\end{align}

\textbf{Step 4: Substitute Lorentz Matrix Elements}

Using $\Lambda^0_{\ 0} = \gamma$, $\Lambda^0_{\ k} = -\gamma v_k/c$,
$\Lambda^i_{\ 0} = -\gamma v^i/c$, and \eqref{eq:lorentz_spatial}:

\textbf{First term:}
\begin{equation}
\Lambda^0_{\ 0} \Lambda^i_{\ j} E^j = \gamma \left[\delta^i_j + (\gamma-1)\frac{v^i v_j}{v^2}\right] E^j
= \gamma E^i + \gamma(\gamma-1)\frac{v^i (\bm{v} \cdot \bm{E})}{v^2}
\end{equation}

\textbf{Second term:}
\begin{equation}
-\Lambda^0_{\ k} \Lambda^i_{\ 0} E^k = -\left(-\frac{\gamma v_k}{c}\right)
\left(-\frac{\gamma v^i}{c}\right) E^k
= -\frac{\gamma^2 v^i (\bm{v} \cdot \bm{E})}{c^2}
\end{equation}

\textbf{Third term:}
\begin{align}
-c \Lambda^0_{\ k} \Lambda^i_{\ l} \epsilon^{klm} B_m
&= -c \left(-\frac{\gamma v_k}{c}\right)
\left[\delta^i_l + (\gamma-1)\frac{v^i v_l}{v^2}\right] \epsilon^{klm} B_m \\
&= \gamma v_k \epsilon^{kim} B_m + \gamma(\gamma-1)\frac{v^i v_l v_k}{v^2} \epsilon^{klm} B_m
\end{align}

The second part vanishes because $v_l v_k \epsilon^{klm} = 0$ (contraction of
symmetric and antisymmetric tensors). The first part gives:
\begin{equation}
\gamma v_k \epsilon^{kim} B_m = \gamma (\bm{v} \times \bm{B})^i
\end{equation}

\textbf{Step 5: Combine All Terms}

\begin{align}
E'^i &= \gamma E^i + \gamma(\gamma-1)\frac{v^i (\bm{v} \cdot \bm{E})}{v^2}
- \frac{\gamma^2 v^i (\bm{v} \cdot \bm{E})}{c^2}
+ \gamma (\bm{v} \times \bm{B})^i
\end{align}

Rearranging:
\begin{align}
E'^i &= \gamma \left[E^i + (\bm{v} \times \bm{B})^i\right]
+ v^i (\bm{v} \cdot \bm{E}) \left[\frac{\gamma(\gamma-1)}{v^2} - \frac{\gamma^2}{c^2}\right]
\end{align}

Simplifying the coefficient using $v^2/c^2 = 1 - 1/\gamma^2$:
\begin{align}
\frac{\gamma(\gamma-1)}{v^2} - \frac{\gamma^2}{c^2}
&= \frac{\gamma(\gamma-1)}{v^2} - \frac{\gamma^2 v^2}{c^2 v^2} \\
&= \frac{\gamma}{v^2}\left[(\gamma-1) - \gamma \cdot \frac{v^2}{c^2}\right] \\
&= \frac{\gamma}{v^2}\left[\gamma - 1 - \gamma + \frac{1}{\gamma}\right] \\
&= \frac{\gamma}{v^2}\left[\frac{1}{\gamma} - 1\right] \\
&= \frac{1 - \gamma}{v^2} = -\frac{\gamma - 1}{v^2}
\end{align}

\textbf{Step 6: Final Result}

\begin{equation}
\boxed{
\bm{E}' = \gamma(\bm{E} + \bm{v} \times \bm{B}) - \frac{\gamma - 1}{v^2}(\bm{v} \cdot \bm{E})\bm{v}
}
\label{eq:E_transform}
\end{equation}

This can also be written as:
\begin{equation}
\bm{E}' = \gamma(\bm{E} + \bm{v} \times \bm{B}) - \frac{\gamma^2}{\gamma + 1}
\frac{(\bm{v} \cdot \bm{E})\bm{v}}{c^2}
\label{eq:E_transform_alt}
\end{equation}

using the identity $(\gamma-1)/v^2 = \gamma^2/[(\gamma+1)c^2]$.

\textbf{Step 7: Parallel and Perpendicular Decomposition}

Decompose $\bm{E} = \bm{E}_\parallel + \bm{E}_\perp$ where
$\bm{E}_\parallel = (\bm{E} \cdot \hat{\bm{v}})\hat{\bm{v}}$ is parallel to $\bm{v}$
and $\bm{E}_\perp = \bm{E} - \bm{E}_\parallel$ is perpendicular.

Since $\bm{v} \times \bm{B}$ is perpendicular to $\bm{v}$:
\begin{align}
\bm{E}'_\parallel &= \gamma \bm{E}_\parallel - (\gamma - 1)\bm{E}_\parallel = \bm{E}_\parallel
\label{eq:E_parallel_transform} \\
\bm{E}'_\perp &= \gamma(\bm{E}_\perp + \bm{v} \times \bm{B})
\label{eq:E_perp_transform}
\end{align}

\textbf{Key result:} The parallel component is unchanged, while the perpendicular
component mixes with the magnetic field.

\textbf{Step 8: Non-Relativistic Limit}

In the limit $v \ll c$, we have $\gamma \approx 1$ and $(\gamma-1)/v^2 \approx 1/(2c^2)$:
\begin{equation}
\bm{E}' \approx \bm{E} + \bm{v} \times \bm{B} + \mathcal{O}(v^2/c^2)
\label{eq:E_transform_nonrel}
\end{equation}

This confirms that $\bm{E} + \bm{v} \times \bm{B}$ is the electric field in the
fluid rest frame (to leading order in $v/c$).

\textbf{Step 9: Magnetic Field Transformation (for completeness)}

By similar derivation (or using duality), the magnetic field transforms as:
\begin{equation}
\boxed{
\bm{B}' = \gamma\left(\bm{B} - \frac{\bm{v} \times \bm{E}}{c^2}\right)
- \frac{\gamma - 1}{v^2}(\bm{v} \cdot \bm{B})\bm{v}
}
\label{eq:B_transform}
\end{equation}

Or in parallel/perpendicular form:
\begin{align}
\bm{B}'_\parallel &= \bm{B}_\parallel \\
\bm{B}'_\perp &= \gamma\left(\bm{B}_\perp - \frac{\bm{v} \times \bm{E}}{c^2}\right)
\end{align}

\subsubsection{The Perfect Conductivity Limit}

Now consider the limit $\sigma \to \infty$ (perfect conductivity). From Ohm's law
\eqref{eq:ohms_law_moving}:
\begin{equation}
\bm{E} + \bm{v} \times \bm{B} = \frac{\bm{J}}{\sigma}
\label{eq:ohms_law_rearranged}
\end{equation}

For any \textit{finite} current density $\bm{J}$, as $\sigma \to \infty$:
\begin{equation}
\bm{E} + \bm{v} \times \bm{B} \to 0
\label{eq:perfect_conductor_limit}
\end{equation}

This means the electric field in the fluid rest frame vanishes:
\begin{equation}
\boxed{\bm{E}' = 0 \quad \Leftrightarrow \quad \bm{E} = -\bm{v} \times \bm{B}}
\label{eq:ideal_mhd_condition}
\end{equation}

\textbf{Physical interpretation:} In a perfect conductor, any electric field would
instantly drive infinite currents that would immediately cancel the field. Therefore,
no electric field can exist in the rest frame of a perfectly conducting fluid.

\subsubsection{The Frozen-In Flux Theorem}

The condition $\bm{E} = -\bm{v} \times \bm{B}$ has a profound consequence: magnetic
field lines are "frozen" into the fluid and move with it.

\textbf{Proof:} Consider the magnetic flux through a surface $S$ moving with the fluid:
\begin{equation}
\Phi = \int_S \bm{B} \cdot d\bm{A}
\end{equation}

The time derivative of flux through a moving surface is:
\begin{equation}
\frac{d\Phi}{dt} = \int_S \frac{\partial \bm{B}}{\partial t} \cdot d\bm{A}
+ \oint_{\partial S} (\bm{v} \times \bm{B}) \cdot d\bm{l}
\label{eq:flux_derivative}
\end{equation}

Using Faraday's law $\partial \bm{B}/\partial t = -\nabla \times \bm{E}$ and
Stokes' theorem:
\begin{equation}
\frac{d\Phi}{dt} = -\oint_{\partial S} \bm{E} \cdot d\bm{l}
+ \oint_{\partial S} (\bm{v} \times \bm{B}) \cdot d\bm{l}
= -\oint_{\partial S} (\bm{E} + \bm{v} \times \bm{B}) \cdot d\bm{l}
\end{equation}

With the ideal MHD condition $\bm{E} + \bm{v} \times \bm{B} = 0$:
\begin{equation}
\boxed{\frac{d\Phi}{dt} = 0}
\label{eq:frozen_flux}
\end{equation}

\textbf{Conclusion:} The magnetic flux through any surface moving with the fluid
is conserved. Equivalently, magnetic field lines are "frozen" into the fluid---they
move, stretch, and compress with the fluid motion.

\subsubsection{Covariant Formulation}

In fully relativistic (covariant) form, the ideal MHD condition is elegantly expressed as:
\begin{equation}
u_\mu F^{\mu\nu} = 0
\label{eq:ideal_mhd_covariant}
\end{equation}

where $u^\mu$ is the fluid four-velocity and $F^{\mu\nu}$ is the electromagnetic
field tensor. This equation states that the electric field vanishes in the fluid
rest frame, which is precisely the ideal MHD condition.

Expanding in components:
\begin{align}
u_\mu F^{\mu 0} &= \gamma(-E^x v_x - E^y v_y - E^z v_z)/c = -\gamma \bm{v} \cdot \bm{E}/c = 0 \\
u_\mu F^{\mu i} &= \gamma(E^i/c + \epsilon^{ijk} v_j B_k) = \gamma(E^i + (\bm{v} \times \bm{B})^i)/c = 0
\end{align}

This recovers $\bm{E} = -\bm{v} \times \bm{B}$ and additionally $\bm{v} \cdot \bm{E} = 0$
(which is automatically satisfied by the cross product).

\subsubsection{When Does Ideal MHD Apply?}

The ideal MHD approximation is valid when:
\begin{enumerate}
    \item \textbf{High conductivity:} The magnetic Reynolds number $R_m = \mu_0 \sigma L V \gg 1$,
        where $L$ and $V$ are characteristic length and velocity scales.
    \item \textbf{Collisional plasma:} The mean free path is much smaller than the
        system size, so the fluid description is valid.
    \item \textbf{Large scales:} We are not resolving resistive layers, reconnection
        regions, or other structures where resistivity matters.
\end{enumerate}

In astrophysical plasmas (stellar interiors, accretion disks, interstellar medium),
$R_m$ can be enormous ($10^{10}$ or larger), making ideal MHD an excellent approximation
on large scales.

%=============================================================================
\section{Stress-Energy Tensor: Derivation from First Principles}
%=============================================================================

\subsection{Perfect Fluid Stress-Energy Tensor}

For a perfect fluid (no viscosity or heat conduction), the stress-energy tensor
in the rest frame is:
\begin{equation}
T^{\mu\nu}_\text{rest} = \begin{pmatrix}
e & 0 & 0 & 0 \\
0 & p & 0 & 0 \\
0 & 0 & p & 0 \\
0 & 0 & 0 & p
\end{pmatrix}
\label{eq:Tmunu_rest}
\end{equation}

where $e = \rho_0 c^2 + \rho_0 \epsilon$ is the energy density and $p$ is the pressure.

To obtain the tensor in an arbitrary frame, we use the covariant form:
\begin{equation}
T^{\mu\nu}_\text{fluid} = (e + p)\frac{u^\mu u^\nu}{c^2} + p \eta^{\mu\nu}
= \rho_0 h \frac{u^\mu u^\nu}{c^2} + p \eta^{\mu\nu}
\label{eq:Tmunu_fluid}
\end{equation}

\subsubsection{Verification of Components}

In the lab frame with fluid velocity $\bm{v}$:
\begin{align}
T^{00} &= \rho_0 h \gamma^2 - p = (e + p)\gamma^2 - p \label{eq:T00_fluid} \\
T^{0i} &= \rho_0 h \gamma^2 \frac{v^i}{c} \label{eq:T0i_fluid} \\
T^{ij} &= \rho_0 h \gamma^2 \frac{v^i v^j}{c^2} + p\delta^{ij} \label{eq:Tij_fluid}
\end{align}

\subsection{Electromagnetic Stress-Energy Tensor}

The electromagnetic stress-energy tensor is:
\begin{equation}
T^{\mu\nu}_\text{EM} = \frac{1}{\mu_0}\left(F^{\mu\alpha}F^\nu_{\ \alpha}
- \frac{1}{4}\eta^{\mu\nu}F^{\alpha\beta}F_{\alpha\beta}\right)
\label{eq:Tmunu_EM}
\end{equation}

In terms of $\bm{E}$ and $\bm{B}$:
\begin{align}
T^{00}_\text{EM} &= \frac{1}{2\mu_0}\left(\frac{E^2}{c^2} + B^2\right) = u_\text{EM} \label{eq:T00_EM} \\
T^{0i}_\text{EM} &= \frac{1}{\mu_0 c}(\bm{E} \times \bm{B})^i = \frac{S^i}{c^2} \label{eq:T0i_EM} \\
T^{ij}_\text{EM} &= \frac{1}{\mu_0}\left[-E^i E^j/c^2 - B^i B^j
+ \frac{1}{2}\delta^{ij}\left(\frac{E^2}{c^2} + B^2\right)\right] \label{eq:Tij_EM}
\end{align}

where $u_\text{EM}$ is the electromagnetic energy density and $\bm{S}$ is the Poynting vector.

\subsection{MHD Stress-Energy Tensor}

In ideal MHD with $\bm{E} = -\bm{v} \times \bm{B}$, we combine the fluid and
electromagnetic contributions. Using the magnetic four-vector $b^\mu$, the
total stress-energy tensor is:
\begin{equation}
\boxed{
T^{\mu\nu} = \left(\rho_0 h + \frac{b^2}{\mu_0}\right)\frac{u^\mu u^\nu}{c^2}
+ \left(p + \frac{b^2}{2\mu_0}\right)\eta^{\mu\nu}
- \frac{b^\mu b^\nu}{\mu_0}
}
\label{eq:Tmunu_MHD}
\end{equation}

\subsubsection{Derivation of MHD Stress-Energy Tensor}

Starting from the EM stress-energy tensor \eqref{eq:Tmunu_EM} and using the ideal
MHD condition $\bm{E} = -\bm{v} \times \bm{B}$:

\textbf{Step 1:} Express $E^2$ and $\bm{E} \times \bm{B}$ in terms of $\bm{v}$ and $\bm{B}$:
\begin{align}
E^2 &= |\bm{v} \times \bm{B}|^2 = v^2 B^2 - (\bm{v} \cdot \bm{B})^2 \\
\bm{E} \times \bm{B} &= -(\bm{v} \times \bm{B}) \times \bm{B}
= \bm{v}B^2 - \bm{B}(\bm{v} \cdot \bm{B})
\end{align}

\textbf{Step 2:} The EM energy density becomes:
\begin{equation}
u_\text{EM} = \frac{1}{2\mu_0}\left[\frac{v^2 B^2 - (\bm{v} \cdot \bm{B})^2}{c^2} + B^2\right]
= \frac{B^2}{2\mu_0}\left[1 + \frac{v^2}{c^2} - \frac{(\bm{v} \cdot \bm{B})^2}{c^2 B^2}\right]
\end{equation}

\textbf{Step 3:} Using $b^2 = B^2/\gamma^2 + (\bm{v} \cdot \bm{B})^2/c^2$, after algebraic
manipulation:
\begin{equation}
u_\text{EM} = \frac{b^2}{2\mu_0} + \frac{\gamma^2}{c^2}\left(\frac{b^2}{\mu_0}\right)v^2
- \frac{(b^0)^2}{\mu_0}
\end{equation}

\textbf{Step 4:} Combining with the fluid contribution and collecting terms proportional
to $u^\mu u^\nu$, $\eta^{\mu\nu}$, and $b^\mu b^\nu$, we obtain \eqref{eq:Tmunu_MHD}.

\subsection{Components in Lab Frame}

Using Gaussian units (where factors of $\mu_0$ become $4\pi$) and $c = 1$:
\begin{align}
T^{00} &= \rho_0 h \gamma^2 + \frac{b^2}{4\pi} - p + \frac{B^2 - (b^0)^2}{8\pi}
\label{eq:T00_lab} \\
T^{0i} &= \left(\rho_0 h \gamma^2 + \frac{b^2}{4\pi}\right)v^i - \frac{b^0 b^i}{4\pi}
\label{eq:T0i_lab} \\
T^{ij} &= \left(\rho_0 h \gamma^2 + \frac{b^2}{4\pi}\right)v^i v^j
+ \left(p + \frac{b^2}{8\pi}\right)\delta^{ij} - \frac{b^i b^j}{4\pi}
\label{eq:Tij_lab}
\end{align}

Define the total enthalpy density including magnetic contribution:
\begin{equation}
w^* = \rho_0 h + \frac{b^2}{4\pi}
\label{eq:total_enthalpy}
\end{equation}

And total pressure:
\begin{equation}
p_\text{tot} = p + p_\text{mag} = p + \frac{b^2}{8\pi}
\label{eq:total_pressure_def}
\end{equation}

%=============================================================================
\section{Conservation Laws of SRMHD}
%=============================================================================

\subsection{Baryon Number Conservation}

The baryon number current four-vector is:
\begin{equation}
J^\mu_B = n u^\mu = n \gamma (c, \bm{v})
\label{eq:baryon_current}
\end{equation}

Conservation of baryon number:
\begin{equation}
\partial_\mu J^\mu_B = 0 \quad \Rightarrow \quad
\frac{\partial (n\gamma)}{\partial t} + \nabla \cdot (n\gamma \bm{v}) = 0
\label{eq:baryon_conservation}
\end{equation}

Defining the lab-frame baryon density $D = n\gamma$:
\begin{equation}
\frac{\partial D}{\partial t} + \nabla \cdot (D\bm{v}) = 0
\label{eq:D_conservation}
\end{equation}

\subsection{Energy-Momentum Conservation}

The divergence of the stress-energy tensor vanishes:
\begin{equation}
\partial_\nu T^{\mu\nu} = 0
\label{eq:Tmunu_conservation}
\end{equation}

\subsubsection{Energy Conservation (mu = 0)}

\begin{equation}
\frac{\partial T^{00}}{\partial t} + \partial_i T^{0i} = 0
\label{eq:energy_cons_1}
\end{equation}

Define conserved energy density:
\begin{equation}
\tau = T^{00} - D c^2 = w^* \gamma^2 - p_\text{tot} - \frac{(b^0)^2}{4\pi} - D c^2
\label{eq:tau_def}
\end{equation}

This represents the total energy density minus rest-mass energy.

\subsubsection{Momentum Conservation (mu = i)}

\begin{equation}
\frac{\partial T^{0i}}{\partial t} + \partial_j T^{ij} = 0
\label{eq:momentum_cons_1}
\end{equation}

Define conserved momentum density:
\begin{equation}
S^i = T^{0i}/c = w^* \gamma^2 v^i/c - b^0 b^i/(4\pi c)
\label{eq:Si_def}
\end{equation}

\subsection{Magnetic Field Evolution}

From $\partial_\nu {}^*F^{\mu\nu} = 0$ and the ideal MHD condition:
\begin{equation}
\frac{\partial \bm{B}}{\partial t} - \nabla \times (\bm{v} \times \bm{B}) = 0
\label{eq:induction_eq}
\end{equation}

This can be rewritten as:
\begin{equation}
\frac{\partial \bm{B}}{\partial t} + \nabla \cdot (\bm{v}\bm{B}^T - \bm{B}\bm{v}^T) = 0
\label{eq:induction_conservative}
\end{equation}

\subsection{Conservative Form Summary}

The SRMHD equations in conservative form:
\begin{equation}
\frac{\partial \bm{U}}{\partial t} + \nabla \cdot \bm{F} = 0
\label{eq:conservative_form}
\end{equation}

with conserved variables:
\begin{equation}
\bm{U} = \begin{pmatrix} D \\ S^x \\ S^y \\ S^z \\ \tau \\ B^x \\ B^y \\ B^z \end{pmatrix}
\label{eq:conserved_variables}
\end{equation}

and flux (in $x$-direction):
\begin{equation}
\bm{F}^x = \begin{pmatrix}
D v^x \\
S^x v^x + p_\text{tot} - b^x b^x/(4\pi) \\
S^y v^x - b^x b^y/(4\pi) \\
S^z v^x - b^x b^z/(4\pi) \\
S^x - D v^x \\
0 \\
B^y v^x - B^x v^y \\
B^z v^x - B^x v^z
\end{pmatrix}
\label{eq:flux_x}
\end{equation}

%=============================================================================
\section{Primitive Variable Recovery: Complete Derivation}
%=============================================================================

This section provides the \textbf{complete derivation} of primitive variable recovery
from conserved variables, which is one of the most challenging aspects of SRMHD.

\subsection{The Inversion Problem}

Given the conserved variables $(D, \bm{S}, \tau, \bm{B})$, we must recover the
primitive variables $(\rho_0, \bm{v}, p, \bm{B})$. The magnetic field $\bm{B}$
is both conserved and primitive, but the others require nonlinear inversion.

\subsubsection{Why This Is Difficult}

The relations between conserved and primitive variables involve products and
ratios with the Lorentz factor $\gamma = (1 - v^2)^{-1/2}$:
\begin{align}
D &= \rho_0 \gamma \label{eq:D_def} \\
S^i &= w^* \gamma^2 v^i - \frac{b^0 b^i}{4\pi} \label{eq:S_def_full} \\
\tau &= w^* \gamma^2 - p_\text{tot} - \frac{(b^0)^2}{4\pi} - D \label{eq:tau_def_full}
\end{align}

where:
\begin{align}
w^* &= \rho_0 h + \frac{b^2}{4\pi} \label{eq:wstar} \\
b^0 &= \gamma \bm{v} \cdot \bm{B} \label{eq:b0_def} \\
b^2 &= \frac{B^2}{\gamma^2} + \frac{(\bm{v} \cdot \bm{B})^2}{1} = \frac{B^2 + (b^0)^2}{\gamma^2}
\label{eq:b2_def}
\end{align}

The coupling between $\gamma$, $\bm{v}$, and the thermodynamic variables through
$b^0$ and $b^2$ creates a nonlinear system that cannot be solved analytically.

\subsection{Auxiliary Variables and Reformulation}

Following Noble et al. (2006) and Mignone \& McKinney (2007), we introduce
auxiliary variables to simplify the system.

\subsubsection{Step 1: Define Useful Combinations}

Let:
\begin{align}
Q &= \tau + D \quad \text{(total energy density)} \label{eq:Q_def} \\
\bm{R} &= \bm{S} \quad \text{(momentum density)} \\
B^2 &= \bm{B} \cdot \bm{B} \\
(\bm{S} \cdot \bm{B}) &= S^i B_i
\end{align}

\subsubsection{Step 2: Velocity-Momentum Relation}

From the momentum definition \eqref{eq:S_def_full}:
\begin{equation}
S^i = w^* \gamma^2 v^i - \frac{b^0 b^i}{4\pi}
\label{eq:Si_expanded}
\end{equation}

Using $b^i = B^i/\gamma + \gamma v^i (\bm{v} \cdot \bm{B})$:
\begin{equation}
S^i = w^* \gamma^2 v^i - \frac{b^0}{4\pi}\left(\frac{B^i}{\gamma} + \gamma v^i (\bm{v} \cdot \bm{B})\right)
\end{equation}

Since $b^0 = \gamma (\bm{v} \cdot \bm{B})$:
\begin{equation}
S^i = w^* \gamma^2 v^i - \frac{(\bm{v} \cdot \bm{B}) B^i}{4\pi}
- \frac{\gamma^2 (\bm{v} \cdot \bm{B})^2 v^i}{4\pi}
\label{eq:Si_intermediate}
\end{equation}

Define:
\begin{equation}
\xi = w^* \gamma^2 + \frac{B^2}{4\pi} = \rho_0 h \gamma^2 + \frac{b^2 \gamma^2 + B^2}{4\pi}
\label{eq:xi_def}
\end{equation}

After algebraic manipulation:
\begin{equation}
\bm{S} = \xi \bm{v} - \frac{(\bm{S} \cdot \bm{B})\bm{B}}{4\pi \xi}
\label{eq:S_xi_relation}
\end{equation}

This can be inverted:
\begin{equation}
\bm{v} = \frac{\bm{S} + (\bm{S} \cdot \bm{B})\bm{B}/(4\pi \xi)}{\xi}
= \frac{4\pi \xi \bm{S} + (\bm{S} \cdot \bm{B})\bm{B}}{4\pi \xi^2}
\label{eq:v_from_S}
\end{equation}

\subsubsection{Step 3: Scalar Products}

Taking the dot product of \eqref{eq:S_xi_relation} with $\bm{B}$:
\begin{equation}
\bm{S} \cdot \bm{B} = \xi (\bm{v} \cdot \bm{B}) - \frac{(\bm{S} \cdot \bm{B}) B^2}{4\pi \xi}
\end{equation}

Solving for $\bm{v} \cdot \bm{B}$:
\begin{equation}
\bm{v} \cdot \bm{B} = \frac{\bm{S} \cdot \bm{B}}{\xi}\left(1 + \frac{B^2}{4\pi \xi}\right)
= \frac{(\bm{S} \cdot \bm{B})(4\pi \xi + B^2)}{4\pi \xi^2}
\label{eq:vdotB}
\end{equation}

\subsubsection{Step 4: Velocity Magnitude}

From \eqref{eq:v_from_S}:
\begin{equation}
v^2 = \frac{S^2 + (\bm{S} \cdot \bm{B})^2/(4\pi\xi) + 2(\bm{S} \cdot \bm{B})^2 B^2/(4\pi\xi)^2
+ (\bm{S} \cdot \bm{B})^2 B^2/(4\pi\xi)^2}{\xi^2}
\end{equation}

Simplifying:
\begin{equation}
v^2 = \frac{S^2 \xi^2 + (\bm{S} \cdot \bm{B})^2[\xi + B^2/(4\pi)]^2/\xi^2}{[\xi + B^2/(4\pi)]^2}
\label{eq:v_squared}
\end{equation}

Define $\tilde{\xi} = \xi + B^2/(4\pi)$:
\begin{equation}
v^2 = \frac{S^2}{\tilde{\xi}^2} + \frac{(\bm{S} \cdot \bm{B})^2}{\xi^2 \tilde{\xi}^2}(4\pi \xi + B^2)^2/(4\pi)^2
\label{eq:v2_final}
\end{equation}

\subsection{The One-Dimensional Root-Finding Problem}

\subsubsection{Step 5: Energy Equation Constraint}

From the energy definition:
\begin{equation}
\tau = w^* \gamma^2 - p_\text{tot} - \frac{(b^0)^2}{4\pi} - D
\label{eq:tau_constraint}
\end{equation}

Using $w^* = \rho_0 h + b^2/(4\pi)$ and $p_\text{tot} = p + b^2/(8\pi)$:
\begin{equation}
\tau + D = \rho_0 h \gamma^2 + \frac{b^2 \gamma^2}{4\pi} - p - \frac{b^2}{8\pi} - \frac{(b^0)^2}{4\pi}
\end{equation}

After simplification:
\begin{equation}
Q = \tau + D = \rho_0 h \gamma^2 - p + \frac{b^2(\gamma^2 - 1/2)}{4\pi} - \frac{(b^0)^2}{4\pi}
\label{eq:Q_constraint}
\end{equation}

\subsubsection{Step 6: Introduce Unknown Function}

Let $W = \gamma$ and define the function:
\begin{equation}
Z = \rho_0 h W^2
\label{eq:Z_def}
\end{equation}

Then:
\begin{align}
\xi &= Z + \frac{B^2 + (b^0)^2/W^2}{4\pi} = Z + \frac{B^2}{4\pi} + \frac{(b^0)^2}{4\pi W^2} \\
b^2 &= \frac{B^2 + (b^0)^2}{W^2}
\end{align}

From the equation of state and thermodynamic relations:
\begin{equation}
\rho_0 h = \rho_0 + \frac{\Gamma}{\Gamma-1}p = \frac{D}{W} + \frac{\Gamma}{\Gamma-1}p
\label{eq:rho_h_from_D}
\end{equation}

So:
\begin{equation}
Z = D W + \frac{\Gamma W^2}{\Gamma-1}p
\label{eq:Z_from_p}
\end{equation}

\subsubsection{Step 7: Pressure from Energy Constraint}

From \eqref{eq:Q_constraint}:
\begin{equation}
p = Z - \frac{b^2(W^2 - 1/2)}{4\pi} + \frac{(b^0)^2}{4\pi} - Q
\label{eq:p_from_Z}
\end{equation}

Substituting the expression for $b^2$:
\begin{equation}
p = Z - \frac{(B^2 + (b^0)^2)(W^2 - 1/2)}{4\pi W^2} + \frac{(b^0)^2}{4\pi} - Q
\label{eq:p_from_Z_expanded}
\end{equation}

Simplifying:
\begin{equation}
p = Z - Q - \frac{B^2(W^2 - 1/2)}{4\pi W^2} - \frac{(b^0)^2(W^2 - 1/2)}{4\pi W^2} + \frac{(b^0)^2}{4\pi}
\end{equation}

\begin{equation}
p = Z - Q - \frac{B^2(W^2 - 1/2)}{4\pi W^2} + \frac{(b^0)^2}{4\pi}\left[1 - \frac{W^2 - 1/2}{W^2}\right]
\end{equation}

\begin{equation}
p = Z - Q - \frac{B^2(2W^2 - 1)}{8\pi W^2} + \frac{(b^0)^2}{8\pi W^2}
\label{eq:p_simplified}
\end{equation}

\subsubsection{Step 8: The Root-Finding Function}

We now have two expressions for $Z$:
\begin{enumerate}
    \item From \eqref{eq:Z_from_p}: $Z = DW + \frac{\Gamma W^2}{\Gamma-1}p$
    \item From velocity constraints and definitions
\end{enumerate}

The velocity constraint \eqref{eq:v2_final} gives:
\begin{equation}
v^2 = 1 - \frac{1}{W^2}
\label{eq:v2_from_W}
\end{equation}

Combined with \eqref{eq:v2_final}:
\begin{equation}
1 - \frac{1}{W^2} = \frac{S^2}{\tilde{\xi}^2} + \frac{(\bm{S} \cdot \bm{B})^2(4\pi\xi + B^2)^2}{16\pi^2 \xi^2 \tilde{\xi}^2}
\label{eq:velocity_constraint}
\end{equation}

\textbf{The algorithm:} Given a trial value of $Z$ (or $W$), compute:
\begin{enumerate}
    \item $\xi(Z)$ from the definitions
    \item $v^2(Z)$ from \eqref{eq:velocity_constraint}
    \item $W = (1 - v^2)^{-1/2}$
    \item $(b^0)^2 = W^2 (\bm{v} \cdot \bm{B})^2$ using \eqref{eq:vdotB}
    \item $p(Z, W)$ from \eqref{eq:p_simplified}
    \item $\rho_0 = D/W$
    \item $Z_\text{computed} = \rho_0 h W^2 = DW + \frac{\Gamma W^2}{\Gamma-1}p$
\end{enumerate}

The root-finding problem is:
\begin{equation}
\boxed{f(Z) = Z - Z_\text{computed}(Z) = 0}
\label{eq:root_finding_Z}
\end{equation}

\subsection{Practical Algorithm: 1D Newton-Raphson}

\subsubsection{Newton-Raphson Iteration}

Given initial guess $Z^{(0)}$, iterate:
\begin{equation}
Z^{(n+1)} = Z^{(n)} - \frac{f(Z^{(n)})}{f'(Z^{(n)})}
\label{eq:newton_raphson}
\end{equation}

\subsubsection{Derivative Computation}

The derivative $f'(Z)$ can be computed analytically or numerically. The analytical
form involves:
\begin{equation}
f'(Z) = 1 - \frac{\partial Z_\text{computed}}{\partial Z}
\end{equation}

Each component requires chain-rule differentiation through the constraint equations.

\subsubsection{Initial Guess}

A good initial guess is critical for convergence:
\begin{equation}
Z^{(0)} = D + Q + \frac{B^2}{8\pi}
\label{eq:initial_guess}
\end{equation}

This corresponds to the non-relativistic limit with magnetic contribution.

\subsubsection{Bounds on Z}

Physical constraints require:
\begin{align}
Z &> D \quad \text{(rest mass contribution)} \\
Z &< Z_\text{max} \quad \text{(from } v < 1 \text{)}
\end{align}

\subsection{Alternative: 2D Newton-Raphson}

For better conditioning, one can solve for two unknowns simultaneously.

\subsubsection{Two-Variable System}

Use $(W, Z)$ or $(v^2, p)$ as unknowns. The system becomes:
\begin{align}
f_1(W, Z) &= S^2 - (Z + B^2/(4\pi))^2 v^2(W) + \ldots = 0 \\
f_2(W, Z) &= Q - Z + p(W, Z) + \ldots = 0
\end{align}

\subsubsection{Jacobian}

The Newton-Raphson update:
\begin{equation}
\begin{pmatrix} W^{(n+1)} \\ Z^{(n+1)} \end{pmatrix} =
\begin{pmatrix} W^{(n)} \\ Z^{(n)} \end{pmatrix} -
\bm{J}^{-1} \begin{pmatrix} f_1 \\ f_2 \end{pmatrix}
\end{equation}

where the Jacobian is:
\begin{equation}
\bm{J} = \begin{pmatrix}
\partial f_1/\partial W & \partial f_1/\partial Z \\
\partial f_2/\partial W & \partial f_2/\partial Z
\end{pmatrix}
\label{eq:jacobian}
\end{equation}

\subsection{Special Cases and Robustness}

\subsubsection{Case: Zero Magnetic Field (Pure Hydrodynamics)}

When $\bm{B} = 0$:
\begin{align}
b^0 &= 0, \quad b^2 = 0 \\
\xi &= Z = \rho_0 h W^2 \\
v^2 &= S^2/Z^2
\end{align}

The system reduces to the simpler relativistic hydrodynamics inversion, which
can be solved with a quartic or simpler iteration.

\subsubsection{Case: Perpendicular Momentum (S dot B = 0)}

When momentum is perpendicular to magnetic field:
\begin{equation}
v^2 = \frac{S^2}{\tilde{\xi}^2}
\end{equation}

The system simplifies but still requires iteration due to the coupling through $\tilde{\xi}$.

\subsubsection{Failure Modes}

The inversion can fail when:
\begin{enumerate}
    \item $v^2 \geq 1$: superluminal velocity (unphysical conserved state)
    \item $p < 0$: negative pressure (need floor)
    \item $\rho_0 < 0$: negative density (need floor)
\end{enumerate}

In these cases, apply physical floors and re-solve.

\subsection{Summary: Complete Inversion Algorithm}

\begin{enumerate}
    \item \textbf{Input:} $D$, $\bm{S}$, $\tau$, $\bm{B}$
    \item \textbf{Compute:} $Q = \tau + D$, $S^2 = |\bm{S}|^2$, $B^2 = |\bm{B}|^2$, $\bm{S} \cdot \bm{B}$
    \item \textbf{Initial guess:} $Z^{(0)} = D + Q + B^2/(8\pi)$
    \item \textbf{Iterate} (Newton-Raphson):
    \begin{enumerate}
        \item Compute $\xi$, $\tilde{\xi}$ from $Z$
        \item Compute $v^2$ from constraint \eqref{eq:velocity_constraint}
        \item Compute $W = (1 - v^2)^{-1/2}$
        \item Compute $\bm{v} \cdot \bm{B}$ from \eqref{eq:vdotB}
        \item Compute $(b^0)^2 = W^2 (\bm{v} \cdot \bm{B})^2$
        \item Compute $p$ from \eqref{eq:p_simplified}
        \item Compute $Z_\text{new} = DW + \frac{\Gamma W^2}{\Gamma-1}p$
        \item Update $Z$ via Newton-Raphson
    \end{enumerate}
    \item \textbf{Convergence:} when $|Z_\text{new} - Z| < \epsilon$
    \item \textbf{Output:} $\rho_0 = D/W$, $\bm{v}$ from \eqref{eq:v_from_S}, $p$
\end{enumerate}

%=============================================================================
\section{Wave Structure in SRMHD: Derivation from Dispersion Relation}
%=============================================================================

\subsection{Linearized Equations}

To find wave speeds, linearize about a uniform state $(\rho_0, \bm{v}_0, p_0, \bm{B}_0)$.
Let perturbations be $\sim e^{i(kx - \omega t)}$.

\subsubsection{Perturbation Ansatz}

\begin{equation}
\bm{U} = \bm{U}_0 + \delta\bm{U} \exp[i(k \cdot x - \omega t)]
\label{eq:perturbation}
\end{equation}

where $\delta\bm{U} \ll \bm{U}_0$.

\subsection{Dispersion Relation}

The linearized SRMHD equations in the fluid rest frame give a dispersion relation
that is a polynomial in the phase velocity $\lambda = \omega/k$:
\begin{equation}
\det(\bm{A} - \lambda \bm{I}) = 0
\label{eq:dispersion_general}
\end{equation}

For 1D propagation along $\hat{x}$ with $\bm{B} = (B_x, B_y, 0)$:
\begin{equation}
(\lambda^2 - c_A^2 \cos^2\theta)[(\lambda^2 - c_s^2)(\lambda^2 - c_A^2)
- \lambda^2 c_s^2 c_A^2 \sin^2\theta] = 0
\label{eq:dispersion_1d}
\end{equation}

where $\theta$ is the angle between $\bm{B}$ and the wave vector $\bm{k}$.

\subsection{Characteristic Speeds}

\subsubsection{Alfv\'en Speed}

From the first factor in \eqref{eq:dispersion_1d}:
\begin{equation}
\lambda_A^2 = c_A^2 \cos^2\theta = \frac{b^2 \cos^2\theta}{w^* + b^2}
\label{eq:alfven_speed_derived}
\end{equation}

where:
\begin{equation}
c_A^2 = \frac{b^2}{\rho_0 h + b^2/(4\pi)} = \frac{b^2}{w^*}
\label{eq:alfven_speed_def}
\end{equation}

In the lab frame with velocity $v$, the Alfv\'en wave speeds are:
\begin{equation}
\lambda_A^\pm = \frac{v \cos\alpha \pm c_A \cos\theta}{1 \pm v c_A \cos\theta/c^2}
\label{eq:alfven_lab_frame}
\end{equation}

where $\alpha$ is the angle between $\bm{v}$ and $\bm{k}$.

\subsubsection{Fast and Slow Magnetosonic Speeds}

The second factor gives the magnetosonic speeds. Define:
\begin{equation}
a^2 = c_s^2 + c_A^2 - c_s^2 c_A^2
\label{eq:a_squared}
\end{equation}

Then:
\begin{equation}
\lambda_{f,s}^2 = \frac{a^2 \pm \sqrt{a^4 - 4 c_s^2 c_A^2 \cos^2\theta}}{2}
\label{eq:fast_slow_speeds}
\end{equation}

where $+$ gives fast waves and $-$ gives slow waves.

\subsubsection{Explicit Fast Speed Formula}

\begin{equation}
\boxed{
c_f^2 = \frac{1}{2}\left[(c_s^2 + c_A^2) + \sqrt{(c_s^2 + c_A^2)^2 - 4c_s^2 c_A^2 \cos^2\theta}\right]
}
\label{eq:fast_speed_explicit}
\end{equation}

In the limiting cases:
\begin{itemize}
    \item $\theta = 0$ (parallel): $c_f = \max(c_s, c_A)$, $c_s = \min(c_s, c_A)$
    \item $\theta = \pi/2$ (perpendicular): $c_f = \sqrt{c_s^2 + c_A^2 - c_s^2 c_A^2}$, $c_s = 0$
\end{itemize}

\subsection{Wave Ordering}

The seven characteristic speeds satisfy:
\begin{equation}
-c_f \leq -c_A \leq -c_s \leq 0 \leq c_s \leq c_A \leq c_f
\label{eq:wave_ordering}
\end{equation}

In the lab frame:
\begin{equation}
\lambda_f^- \leq \lambda_A^- \leq \lambda_s^- \leq \lambda_0 \leq \lambda_s^+ \leq \lambda_A^+ \leq \lambda_f^+
\label{eq:lab_frame_ordering}
\end{equation}

where $\lambda_0 = v$ is the entropy/contact wave.

%=============================================================================
\section{Rankine-Hugoniot Conditions for SRMHD Shocks}
%=============================================================================

\subsection{Jump Conditions Derivation}

For a shock moving with speed $V_s$, the Rankine-Hugoniot conditions express
conservation across the discontinuity.

\subsubsection{Frame Transformation}

In the shock rest frame, the flow is steady. Let subscripts $a$ and $b$ denote
upstream (ahead) and downstream (behind) states.

Conservation of baryon number:
\begin{equation}
[\![ \rho_0 \gamma (v - V_s) ]\!] = 0 \quad \Rightarrow \quad
\rho_{0a} \gamma_a (v_a - V_s) = \rho_{0b} \gamma_b (v_b - V_s)
\label{eq:RH_mass}
\end{equation}

Define the mass flux through the shock:
\begin{equation}
J = \rho_0 \gamma (v - V_s) = \text{const}
\label{eq:mass_flux_def}
\end{equation}

\subsubsection{Normal Momentum Conservation}

\begin{equation}
[\![ \rho_0 h \gamma^2 (v - V_s) v_n + p_\text{tot} - \frac{b_n^2}{4\pi} ]\!] = 0
\label{eq:RH_momentum_n}
\end{equation}

Using $J$:
\begin{equation}
J h_a \gamma_a v_{n,a} + p_{\text{tot},a} - \frac{b_{n,a}^2}{4\pi}
= J h_b \gamma_b v_{n,b} + p_{\text{tot},b} - \frac{b_{n,b}^2}{4\pi}
\label{eq:RH_momentum_n_2}
\end{equation}

\subsubsection{Tangential Momentum Conservation}

\begin{equation}
[\![ \rho_0 h \gamma^2 (v - V_s) v_t - \frac{b_n b_t}{4\pi} ]\!] = 0
\label{eq:RH_momentum_t}
\end{equation}

\subsubsection{Energy Conservation}

\begin{equation}
[\![ (\rho_0 h \gamma^2 + b^2/4\pi)(v - V_s) - (b^0/4\pi)(b_n/\gamma) ]\!] = 0
\label{eq:RH_energy}
\end{equation}

\subsubsection{Magnetic Field Conditions}

From $\nabla \cdot \bm{B} = 0$:
\begin{equation}
[\![ B_n ]\!] = 0 \quad \Rightarrow \quad B_{n,a} = B_{n,b}
\label{eq:RH_divB}
\end{equation}

From the induction equation:
\begin{equation}
[\![ (v - V_s) B_t - v_t B_n ]\!] = 0
\label{eq:RH_induction}
\end{equation}

\subsection{Shock Adiabat (Hugoniot)}

Combining the conservation equations, the shock adiabat relates upstream and
downstream thermodynamic states:

For pure gas dynamics ($B = 0$):
\begin{equation}
(h_a - h_b)(p_a + p_b) = (h_a + h_b)(p_a - p_b)/\gamma_\text{eff}
\label{eq:hugoniot_hydro}
\end{equation}

For MHD, the Hugoniot is more complex due to magnetic contributions. The general
form involves:
\begin{equation}
[\![ h ]\!] \cdot [\![ p + b^2/8\pi ]\!] + [\![ h \cdot (p + b^2/8\pi) ]\!] \cdot f(\bm{B}) = 0
\label{eq:hugoniot_mhd}
\end{equation}

\subsection{Fast Shock Solution}

For a fast shock with given upstream state and downstream pressure $p_b$:

\textbf{Step 1:} From the Hugoniot, find $\rho_{0b}$.

\textbf{Step 2:} From mass conservation:
\begin{equation}
J^2 = -\frac{p_{\text{tot},b} - p_{\text{tot},a}}{(h/\rho_0)_b - (h/\rho_0)_a}
\label{eq:J_squared}
\end{equation}

\textbf{Step 3:} Find shock speed:
\begin{equation}
V_s = \frac{v_a + J/(\rho_{0a} \gamma_a^2)}{1 + v_a J/(\rho_{0a} \gamma_a^2 c^2)}
\label{eq:shock_speed_formula}
\end{equation}

\textbf{Step 4:} Find downstream velocity:
\begin{equation}
v_b = \frac{V_s - J/(\rho_{0b}\gamma_b)}{1 - V_s J/(\rho_{0b}\gamma_b c^2)}
\label{eq:downstream_velocity}
\end{equation}

This requires iteration since $\gamma_b$ depends on $v_b$.

%=============================================================================
\section{SPH Discretization}
%=============================================================================

\subsection{Kernel Function and Density Estimation}

Following Kitajima et al. (2025), we use the Gaussian kernel:
\begin{equation}
W(\bm{x}, h) = \left(\frac{1}{\sqrt{\pi} h}\right)^d \exp\left(-\frac{|\bm{x}|^2}{h^2}\right)
\label{eq:kernel}
\end{equation}

where $d$ is the number of spatial dimensions.

\subsubsection{Volume-Based Density Formulation}

Define the particle volume:
\begin{equation}
V_p(\bm{x}) \equiv \left[\sum_j W(\bm{x} - \bm{x}_j, h(\bm{x}))\right]^{-1}
\label{eq:volume}
\end{equation}

The lab-frame number density:
\begin{equation}
D(\bm{x}) = \nu(\bm{x}) / V_p(\bm{x})
\label{eq:number_density}
\end{equation}

where $\nu(\bm{x})$ is the baryon number per particle.

\subsubsection{Variable Smoothing Length}

The smoothing length is determined iteratively:
\begin{equation}
h(\bm{x}) = \eta \left[V_p^*(\bm{x})\right]^{1/d}
\label{eq:smoothing_length}
\end{equation}

where:
\begin{equation}
V_p^*(\bm{x}) = \left[\sum_j W(\bm{x} - \bm{x}_j, C_\text{smooth} h(\bm{x}))\right]^{-1}
\label{eq:volume_star}
\end{equation}

Typical values: $\eta = 1.0$, $C_\text{smooth} = 2.0$.

\subsection{Convolution Formulation}

Physical quantities are defined through convolution:
\begin{equation}
\langle f_i \rangle \equiv \int f(\bm{x}) W(\bm{x} - \bm{x}_i, h(\bm{x})) d^3x
\label{eq:convolution}
\end{equation}

\subsection{Derivation of SRMHD-SPH Equations}

\subsubsection{Momentum Equation}

Starting from the Lagrangian momentum equation:
\begin{equation}
\frac{d\bm{S}}{dt} = -\frac{1}{D} \nabla^\nu T^{\mu\nu}
\end{equation}

Taking the convolution with the kernel:
\begin{align}
\langle \dot{S}_i^\mu \rangle &= \int \frac{d S^\mu(\bm{x})}{dt} W_i(\bm{x}) d^3x \\
&= -\int \frac{1}{D(\bm{x})} \nabla^\nu T^{\mu\nu}(\bm{x}) W_i(\bm{x}) d^3x
\end{align}

Following the derivation in Kitajima et al. (2025):
\begin{align}
\langle \dot{S}_i^\mu \rangle &= -\sum_j \int \frac{T^{\mu\nu}(\bm{x})}{D^2(\bm{x})}
\left[(\nabla_i - \nabla_j) - \frac{h}{2}\nabla h (\nabla_i^2 - \nabla_j^2)\right] \nonumber \\
&\quad \times W(\bm{x}-\bm{x}_i, h(\bm{x})) W(\bm{x}-\bm{x}_j, h(\bm{x})) d^3x
\label{eq:mom_sph_derivation}
\end{align}

\subsubsection{Approximation for Smooth h}

Assuming $\nabla h \approx 0$ (ensured by $C_\text{smooth} > 1$):
\begin{equation}
\langle \nu_i \dot{S}_i^\mu \rangle = -\sum_j (T^{\mu\nu})^* n_\nu V_{ij,\text{interp}}^2
\left[\nabla_i W(\bm{x}_i - \bm{x}_j, \sqrt{2}h_i) - \nabla_j W(\bm{x}_i - \bm{x}_j, \sqrt{2}h_j)\right]
\label{eq:mom_sph}
\end{equation}

where $V_{ij}^2$ is obtained by interpolating $D^{-2}(\bm{x})$.

\subsubsection{Energy Equation}

Similarly, the energy equation becomes:
\begin{equation}
\langle \nu_i \dot{e}_i \rangle = -\sum_j (T^{\mu\nu} v_\nu)^* n_\mu V_{ij,\text{interp}}^2
\left[\nabla_i W(\bm{x}_i - \bm{x}_j, \sqrt{2}h_i) - \nabla_j W(\bm{x}_i - \bm{x}_j, \sqrt{2}h_j)\right]
\label{eq:energy_sph}
\end{equation}

\subsubsection{Induction Equation}

The SPH form of the induction equation:
\begin{equation}
\frac{d}{dt}\left(\frac{B^\mu}{D}\right)_i = \sum_j m_j (B^\nu)^* n_\nu
\left[(v^\mu)^* - \dot{x}_i^{*\mu}\right] F_{ij}
\label{eq:induction_sph}
\end{equation}

where $\dot{\bm{x}}_i^* = \bm{v}_i + \ddot{\bm{x}}_i \Delta t / 2$ is the time-centered velocity.

%=============================================================================
\section{Riemann Solver for SRMHD}
%=============================================================================

\subsection{Decomposition Strategy}

Following Iwasaki \& Inutsuka (2011) and extending to relativistic case, we decompose the
MHD Riemann problem into:

\begin{enumerate}
    \item \textbf{Compressive part}: Riemann problem with total pressure $P_t$ (fast shocks/rarefactions)
    \item \textbf{Alfv\'enic part}: Method of Characteristics for transverse components
\end{enumerate}

\subsection{Relativistic Riemann Problem with Magnetic Pressure}

For the compressive part, we solve a Riemann problem with effective total pressure:
\begin{equation}
P_t = p + \frac{B_\perp^2}{8\pi}
\label{eq:total_pressure}
\end{equation}

\subsubsection{Wave Structure}

The solution contains (from left to right):
\begin{itemize}
    \item Left fast shock or rarefaction (FS/FR)
    \item Contact discontinuity (CD)
    \item Right fast shock or rarefaction (FS/FR)
\end{itemize}

At the contact discontinuity:
\begin{align}
P_t^* &= P_L^* + \frac{(B_{\perp L}^*)^2}{8\pi} = P_R^* + \frac{(B_{\perp R}^*)^2}{8\pi} \\
v_\parallel^* &= v_{\parallel L}^* = v_{\parallel R}^*
\end{align}

\subsubsection{Shock Relations}

For a shock wave, the relativistic Rankine-Hugoniot conditions give:

\textbf{Mass flux:}
\begin{equation}
j^2 = -\frac{P_b - P_a}{H_b/n_b - H_a/n_a}
\label{eq:mass_flux}
\end{equation}

\textbf{Shock speed:}
\begin{equation}
V_s = \frac{D_a^2 v_{x,a} \pm j\sqrt{j^2 + D_a^2(1-v_{x,a}^2)}}{D_a^2 + j^2}
\label{eq:shock_speed}
\end{equation}

where $+$ for right-going shocks and $-$ for left-going shocks.

\textbf{Post-shock velocity:}
\begin{equation}
v_{x,b} = \frac{H_a \gamma_a v_{x,a} + \gamma_s(P_b - P_a)/j_\text{signed}}
{H_a \gamma_a + (P_b - P_a)(\gamma_s v_{x,a}/j_\text{signed} + 1/D_a)}
\label{eq:post_shock_velocity}
\end{equation}

\subsubsection{Rarefaction Relations}

For rarefaction waves, the Riemann invariant is conserved:
\begin{equation}
J_\pm = \frac{1}{2}\ln\left(\frac{1+v_x}{1-v_x}\right) \pm \int \frac{dp}{\rho_0 h c_s}
\label{eq:riemann_invariant}
\end{equation}

The velocity behind the rarefaction:
\begin{equation}
v_{x,b} = \tanh\left[\frac{1}{2}\ln\left(\frac{1+v_{x,a}}{1-v_{x,a}}\right)
\pm \int_{p_a}^{p_b} \frac{dp'}{\rho_0(p') h(p') c_s(p')}\right]
\label{eq:rarefaction_velocity}
\end{equation}

\subsection{Iterative Solution}

The complete Riemann solution is found by root-finding:
\begin{equation}
f(p^*) = v_{x,L}(p^*) - v_{x,R}(p^*) = 0
\label{eq:root_finding}
\end{equation}

where $v_{x,L}(p^*)$ and $v_{x,R}(p^*)$ are computed from the respective wave curves.

\subsection{Tangential Velocity Treatment}

For relativistic flows with tangential velocity $v^t$, the invariant:
\begin{equation}
K = H \gamma v^t
\label{eq:tangent_invariant}
\end{equation}

is conserved across rarefaction waves (Pons et al. 2000).

The tangential velocity in the star region:
\begin{equation}
v_t^* = \frac{K}{\sqrt{H^{*2} + K^2/(1 - v_x^{*2})}}
\label{eq:tangent_star}
\end{equation}

%=============================================================================
\section{Method of Characteristics for Alfv\'en Waves}
%=============================================================================

\subsection{Relativistic Alfv\'en Wave Equations}

The transverse (Alfv\'enic) components satisfy:
\begin{align}
\frac{\partial v_\perp}{\partial t} &= \frac{B_\parallel}{4\pi \rho_0 H \gamma^2} \nabla_\parallel B_\perp \\
\frac{\partial B_\perp}{\partial t} &= B_\parallel \nabla_\parallel v_\perp
\end{align}

These decouple from the compressive modes when $B_\parallel$ is constant.

\subsection{Characteristic Solution}

The Alfv\'en characteristics are:
\begin{equation}
\frac{dx}{dt} = v_\parallel \pm c_A
\label{eq:alfven_characteristic}
\end{equation}

where $c_A$ is the relativistic Alfv\'en speed.

\subsection{Method of Characteristics (MOC)}

At the particle interface, the transverse quantities are determined by:

\begin{equation}
\bm{B}_\perp^* = \frac{\bm{B}_{\perp,L}/\sqrt{\rho_L} + \bm{B}_{\perp,R}/\sqrt{\rho_R}
+ \text{sgn}(B_\parallel)(\bm{v}_{\perp,R} - \bm{v}_{\perp,L})}
{1/\sqrt{\rho_L} + 1/\sqrt{\rho_R}}
\label{eq:moc_B}
\end{equation}

\begin{equation}
\bm{v}_\perp^* = \frac{\bm{v}_{\perp,L}\sqrt{\rho_L} + \bm{v}_{\perp,R}\sqrt{\rho_R}
+ \text{sgn}(B_\parallel)(\bm{B}_{\perp,R} - \bm{B}_{\perp,L})}
{\sqrt{\rho_L} + \sqrt{\rho_R}}
\label{eq:moc_v}
\end{equation}

In the relativistic case, $\rho$ should be replaced by $\rho_0 h \gamma^2$ for proper wave impedance matching:
\begin{equation}
Z_A = \sqrt{\rho_0 h \gamma^2 + b^2/(4\pi)}
\label{eq:alfven_impedance}
\end{equation}

%=============================================================================
\section{Complete SRMHD-GSPH Algorithm}
%=============================================================================

\subsection{Final Discretized Equations}

Combining the Riemann solver for compressive waves and MOC for Alfv\'en waves:

\subsubsection{Momentum Update}

\begin{equation}
\boxed{
\frac{d\bm{S}_i}{dt} = -\sum_j m_j F_{ij}
\left[-\left((P_t)_\text{RP} - P_{B,\parallel}\right)\bm{n} + B_\parallel^* (\bm{B}_\perp)_\text{MOC}\right]
}
\label{eq:final_momentum}
\end{equation}

where:
\begin{itemize}
    \item $F_{ij} = 2 V_{ij}^2 \partial W / \partial s_i$ (kernel gradient factor)
    \item $(P_t)_\text{RP}$ is the total pressure from the Riemann solver
    \item $P_{B,\parallel} = (B_{\parallel,i}^2 + B_{\parallel,j}^2)/4$
    \item $B_\parallel^* = (B_{\parallel,i} + B_{\parallel,j})/2$
    \item $(\bm{B}_\perp)_\text{MOC}$ from Method of Characteristics
\end{itemize}

\subsubsection{Energy Update}

\begin{equation}
\boxed{
\frac{de_i}{dt} = -\sum_j m_j F_{ij}
\left[-\left((P_t)_\text{RP} - P_{B,\parallel}\right)(v_\parallel)_\text{RP}
+ B_\parallel^* (\bm{B}_\perp)_\text{MOC} \cdot (\bm{v}_\perp)_\text{MOC}\right]
}
\label{eq:final_energy}
\end{equation}

\subsubsection{Induction Equation}

\begin{equation}
\boxed{
\frac{d}{dt}\left(\frac{\bm{B}}{D}\right)_i = \sum_j m_j F_{ij} B_\parallel^*
\left[(v_\parallel)_\text{RP} \bm{n} + (\bm{v}_\perp)_\text{MOC} - \dot{\bm{x}}_i^*\right]
}
\label{eq:final_induction}
\end{equation}

\subsection{Maintaining the Divergence-Free Constraint: Complete Treatment}

The constraint $\nabla \cdot \bm{B} = 0$ is fundamental to electromagnetism---there are
no magnetic monopoles. However, numerical discretization generically produces
$\nabla \cdot \bm{B} \neq 0$ errors that can grow and destabilize simulations.
This section presents a comprehensive treatment of methods to handle this constraint.

\subsubsection{Historical Context: From General Relativity to MHD}

The problem of constraint preservation in numerical simulations was first seriously
addressed in \textbf{numerical general relativity}. Einstein's equations contain
constraint equations (the Hamiltonian and momentum constraints) that must be
satisfied on each time slice. Early numerical relativity codes suffered from
constraint-violating modes that grew exponentially, crashing simulations.

The breakthrough came with \textbf{hyperbolic constraint damping}, developed for
GR by Gundlach et al. (2005) and refined by others. The key insight was to
modify the evolution equations to include terms proportional to the constraint
violations, causing these violations to propagate away and decay rather than grow.

This technique was then adapted to MHD by Dedner et al. (2002), where the
constraint $\nabla \cdot \bm{B} = 0$ plays an analogous role to the GR constraints.
The first successful applications were indeed in MHD, where the method proved
remarkably effective. The technique has since been applied back to GR and to
other constrained systems.

\subsubsection{The Divergence Problem in MHD}

The exact induction equation:
\begin{equation}
\frac{\partial \bm{B}}{\partial t} = \nabla \times (\bm{v} \times \bm{B})
\label{eq:exact_induction}
\end{equation}

automatically preserves $\nabla \cdot \bm{B}$:
\begin{equation}
\frac{\partial}{\partial t}(\nabla \cdot \bm{B}) = \nabla \cdot \frac{\partial \bm{B}}{\partial t}
= \nabla \cdot [\nabla \times (\bm{v} \times \bm{B})] = 0
\end{equation}

However, numerical discretization breaks this property. In SPH, the discrete
gradient and curl operators do not exactly satisfy $\nabla \cdot (\nabla \times \bm{A}) = 0$.
Non-zero $\nabla \cdot \bm{B}$ causes:
\begin{enumerate}
    \item \textbf{Spurious parallel forces:} The magnetic tension force
        $(\bm{B} \cdot \nabla)\bm{B}$ acquires unphysical components.
    \item \textbf{Tensile instability:} Particles can clump along field lines.
    \item \textbf{Energy errors:} Non-conservative energy exchange.
\end{enumerate}

\subsubsection{Method 1: Powell Source Terms (8-Wave Formulation)}

Powell et al. (1999) proposed adding source terms proportional to $\nabla \cdot \bm{B}$
to the MHD equations. This is also called the "8-wave" formulation because it
introduces an eighth wave (the divergence wave) in addition to the seven physical
MHD waves.

\textbf{Modified equations:}
\begin{align}
\frac{\partial \rho}{\partial t} + \nabla \cdot (\rho \bm{v}) &= 0 \\
\frac{\partial (\rho \bm{v})}{\partial t} + \nabla \cdot \bm{T} &= -\bm{B}(\nabla \cdot \bm{B})
\label{eq:powell_momentum_pde} \\
\frac{\partial E}{\partial t} + \nabla \cdot [(E + p_\text{tot})\bm{v} - \bm{B}(\bm{v} \cdot \bm{B})]
&= -(\bm{v} \cdot \bm{B})(\nabla \cdot \bm{B}) \label{eq:powell_energy_pde} \\
\frac{\partial \bm{B}}{\partial t} + \nabla \cdot (\bm{v}\bm{B}^T - \bm{B}\bm{v}^T)
&= -\bm{v}(\nabla \cdot \bm{B}) \label{eq:powell_induction_pde}
\end{align}

\textbf{Physical interpretation:}
\begin{itemize}
    \item The source terms transport $\nabla \cdot \bm{B}$ errors with the fluid velocity $\bm{v}$.
    \item Errors are advected out of the computational domain.
    \item Does \textit{not} reduce $\nabla \cdot \bm{B}$, only prevents its growth.
\end{itemize}

\textbf{SPH discretization:}

\begin{equation}
\frac{d\bm{S}_i}{dt} \rightarrow \frac{d\bm{S}_i}{dt} - \bm{B}_i \sum_j m_j B_\parallel^* F_{ij}
\label{eq:powell_momentum}
\end{equation}

\begin{equation}
\frac{de_i}{dt} \rightarrow \frac{de_i}{dt} - (\bm{B}_i \cdot \dot{\bm{x}}_i^*) \sum_j m_j B_\parallel^* F_{ij}
\label{eq:powell_energy}
\end{equation}

\textbf{Advantages:}
\begin{itemize}
    \item Simple to implement
    \item No additional evolution equations
    \item Works well for outflow boundaries
\end{itemize}

\textbf{Disadvantages:}
\begin{itemize}
    \item Breaks strict conservation of momentum and energy
    \item Does not actively reduce existing $\nabla \cdot \bm{B}$ errors
    \item Can accumulate errors in periodic domains
\end{itemize}

\subsubsection{Method 2: Hyperbolic Divergence Cleaning (Dedner et al.)}

Dedner et al. (2002) introduced a more sophisticated approach that actively
damps $\nabla \cdot \bm{B}$ errors. This method, inspired by constraint damping
in numerical relativity, introduces an auxiliary scalar field $\psi$ coupled
to the magnetic field evolution.

\textbf{Generalized Lagrangian Multiplier (GLM) formulation:}

Replace the induction equation with:
\begin{align}
\frac{\partial \bm{B}}{\partial t} + \nabla \cdot (\bm{v}\bm{B}^T - \bm{B}\bm{v}^T) + \nabla\psi &= 0
\label{eq:glm_induction} \\
\frac{\partial \psi}{\partial t} + c_h^2 \nabla \cdot \bm{B} &= -\frac{c_h^2}{c_p^2}\psi
\label{eq:glm_psi}
\end{align}

where:
\begin{itemize}
    \item $\psi$ is the divergence cleaning potential (scalar field)
    \item $c_h$ is the hyperbolic cleaning speed (wave propagation)
    \item $c_p$ is the parabolic damping rate
\end{itemize}

\textbf{Derivation of the cleaning mechanism:}

Taking the divergence of \eqref{eq:glm_induction}:
\begin{equation}
\frac{\partial (\nabla \cdot \bm{B})}{\partial t} + \nabla^2 \psi = 0
\end{equation}

Taking the time derivative and substituting \eqref{eq:glm_psi}:
\begin{equation}
\frac{\partial^2 (\nabla \cdot \bm{B})}{\partial t^2}
- c_h^2 \nabla^2 (\nabla \cdot \bm{B})
+ \frac{c_h^2}{c_p^2} \frac{\partial (\nabla \cdot \bm{B})}{\partial t} = 0
\label{eq:telegraph}
\end{equation}

This is the \textbf{telegraph equation}---a damped wave equation! It shows that
$\nabla \cdot \bm{B}$ errors:
\begin{enumerate}
    \item Propagate as waves with speed $c_h$ (hyperbolic behavior)
    \item Decay exponentially with rate $c_h^2/(2c_p^2)$ (parabolic damping)
\end{enumerate}

\textbf{Parameter choices:}

Typical choices following Tricco \& Price (2012):
\begin{align}
c_h &= \max_i(v_{\text{fast},i}) \quad \text{(fastest wave speed)} \\
c_p &= \frac{c_h h}{\sigma} \quad \text{where } \sigma \approx 0.2\text{--}1.0
\end{align}

The ratio $c_h/c_p$ controls the balance between propagation and damping.

\textbf{SPH implementation:}

The $\psi$ field is evolved as an additional variable per particle:
\begin{equation}
\frac{d\psi_i}{dt} = -c_h^2 (\nabla \cdot \bm{B})_i - \frac{c_h^2}{c_p^2}\psi_i
\label{eq:psi_evolution_sph}
\end{equation}

The magnetic field evolution gains a gradient term:
\begin{equation}
\frac{d\bm{B}_i}{dt} = \ldots - (\nabla\psi)_i
\label{eq:B_psi_correction}
\end{equation}

\textbf{Advantages:}
\begin{itemize}
    \item Actively reduces $\nabla \cdot \bm{B}$ errors
    \item Errors propagate away at finite speed $c_h$
    \item Works in periodic domains
    \item Maintains conservation to machine precision (with careful implementation)
\end{itemize}

\textbf{Disadvantages:}
\begin{itemize}
    \item Requires evolving an additional field $\psi$
    \item Extra computational cost and memory
    \item Requires careful parameter tuning ($c_h$, $c_p$)
\end{itemize}

\subsubsection{Method 3: Constrained Transport (CT)}

Constrained Transport, developed by Evans \& Hawley (1988), is the gold standard
for maintaining $\nabla \cdot \bm{B} = 0$ exactly (to machine precision) in
grid-based methods. While not directly applicable to SPH, understanding CT
provides insight into the constraint preservation problem.

\textbf{Key idea:} Store magnetic field components at \textit{face centers}
of a staggered grid, and electric field (EMF) at \textit{edge centers}.

For a 2D grid cell with corners at $(i,j)$, $(i+1,j)$, $(i,j+1)$, $(i+1,j+1)$:
\begin{itemize}
    \item $B_x$ stored at $(i+1/2, j)$ and $(i+1/2, j+1)$ (x-faces)
    \item $B_y$ stored at $(i, j+1/2)$ and $(i+1, j+1/2)$ (y-faces)
    \item $E_z$ stored at $(i+1/2, j+1/2)$ (z-edge, center of cell)
\end{itemize}

\textbf{Discrete divergence:}
\begin{equation}
(\nabla \cdot \bm{B})_{i,j} = \frac{B_x|_{i+1/2,j} - B_x|_{i-1/2,j}}{\Delta x}
+ \frac{B_y|_{i,j+1/2} - B_y|_{i,j-1/2}}{\Delta y}
\label{eq:ct_divB}
\end{equation}

\textbf{Induction equation update:}
\begin{align}
\frac{\partial B_x|_{i+1/2,j}}{\partial t} &= -\frac{E_z|_{i+1/2,j+1/2} - E_z|_{i+1/2,j-1/2}}{\Delta y}
\label{eq:ct_Bx} \\
\frac{\partial B_y|_{i,j+1/2}}{\partial t} &= +\frac{E_z|_{i+1/2,j+1/2} - E_z|_{i-1/2,j+1/2}}{\Delta x}
\label{eq:ct_By}
\end{align}

\textbf{Why it preserves $\nabla \cdot \bm{B} = 0$:}

Taking the discrete time derivative of \eqref{eq:ct_divB} and substituting
\eqref{eq:ct_Bx}--\eqref{eq:ct_By}, all $E_z$ terms cancel in pairs:
\begin{equation}
\frac{\partial}{\partial t}(\nabla \cdot \bm{B})_{i,j} = 0 \quad \text{(exactly)}
\end{equation}

This is a discrete analogue of $\nabla \cdot (\nabla \times \bm{E}) = 0$.

\textbf{CT for SPH (Vector Potential Method):}

Price \& Monaghan (2004) and Price (2012) developed an SPH analogue using the
vector potential $\bm{A}$ where $\bm{B} = \nabla \times \bm{A}$:
\begin{equation}
\frac{d\bm{A}_i}{dt} = \bm{v}_i \times \bm{B}_i + \nabla\phi_i
\label{eq:vector_potential_sph}
\end{equation}

where $\phi$ is a gauge freedom. Then:
\begin{equation}
\bm{B}_i = \sum_j \frac{m_j}{\rho_j} \bm{A}_j \times \nabla_i W_{ij}
\label{eq:B_from_A}
\end{equation}

This automatically satisfies $\nabla \cdot \bm{B} = 0$ because
$\nabla \cdot (\nabla \times \bm{A}) = 0$ is an identity.

\textbf{Advantages of vector potential:}
\begin{itemize}
    \item Exactly divergence-free by construction
    \item No cleaning needed
    \item Preserves magnetic helicity
\end{itemize}

\textbf{Disadvantages:}
\begin{itemize}
    \item Gauge freedom requires careful handling
    \item Can develop spurious modes
    \item Less intuitive than evolving $\bm{B}$ directly
    \item Boundary conditions more complex
\end{itemize}

\subsubsection{Comparison of Divergence Cleaning Methods}

\begin{table}[h]
\centering
\begin{tabular}{|l|c|c|c|c|}
\hline
\textbf{Method} & \textbf{$\nabla \cdot B = 0$} & \textbf{Conservation} & \textbf{Cost} & \textbf{SPH} \\
\hline
Powell (8-wave) & Advected & Broken & Low & Easy \\
\hline
Dedner (GLM-$\psi$) & Damped & Good & Medium & Yes \\
\hline
CT (grid) & Exact & Exact & Low & No \\
\hline
Vector potential & Exact & Good & Medium & Yes \\
\hline
Projection & Exact & Broken & High & Difficult \\
\hline
\end{tabular}
\caption{Comparison of divergence-free maintenance methods}
\label{tab:divB_methods}
\end{table}

\subsubsection{Recommendation for SRMHD-GSPH}

For our SRMHD-GSPH implementation, we recommend a \textbf{hybrid approach}:

\begin{enumerate}
    \item \textbf{Primary:} Hyperbolic divergence cleaning (Dedner/GLM)
        \begin{itemize}
            \item Actively reduces errors
            \item Works with Riemann solver framework
            \item Well-tested in relativistic MHD codes
        \end{itemize}

    \item \textbf{Supplementary:} Powell source terms
        \begin{itemize}
            \item Cheap insurance against residual errors
            \item Helps with boundary effects
        \end{itemize}
\end{enumerate}

The relativistic generalization of the GLM system:
\begin{align}
\frac{\partial \bm{B}}{\partial t} + \nabla \cdot (\bm{v}\bm{B}^T - \bm{B}\bm{v}^T) + \nabla\psi &= 0 \\
\frac{\partial \psi}{\partial t} + c_h^2 \nabla \cdot \bm{B} + \frac{c_h^2}{c_p^2}\psi &= 0
\end{align}

where $c_h$ should be the relativistic fast magnetosonic speed to ensure proper
CFL stability.

\subsection{Time Integration}

Use explicit time-stepping (e.g., Euler or Runge-Kutta):

\begin{align}
(\nu_i \bm{S}_i)^{n+1} &= (\nu_i \bm{S}_i)^n + (\nu_i \dot{\bm{S}}_i) \Delta t \\
(\nu_i e_i)^{n+1} &= (\nu_i e_i)^n + (\nu_i \dot{e}_i) \Delta t \\
\bm{x}_i^{n+1} &= \bm{x}_i^n + \bm{v}_i \Delta t \\
(\bm{B}/D)_i^{n+1} &= (\bm{B}/D)_i^n + (d/dt)(\bm{B}/D)_i \Delta t
\end{align}

\subsection{CFL Condition}

The timestep is limited by:
\begin{equation}
\Delta t = C_\text{CFL} \min_i \left[\frac{h_i}{c_{f,i}}\right]
\label{eq:cfl}
\end{equation}

where $c_f$ is the relativistic fast magnetosonic speed and $C_\text{CFL} \approx 0.3$.

%=============================================================================
\section{Implementation Notes}
%=============================================================================

\subsection{Numerical Stability}

\begin{enumerate}
    \item \textbf{Monotonicity constraint}: Apply limiters to gradients in MUSCL reconstruction
    \item \textbf{Positivity preservation}: Floor pressure and density to small positive values
    \item \textbf{Subluminal velocities}: Ensure $|\bm{v}| < c$ after each update
\end{enumerate}

\subsection{Shock Detection}

Switch to first-order (no reconstruction) when:
\begin{equation}
C_\text{shock} \bm{e}_{ij} \cdot (\bm{v}_i - \bm{v}_j) > \min(c_{s,i}, c_{s,j})
\quad \text{or} \quad
\left|\log_{10}(p_i/p_j)\right| > C_\text{cd}
\label{eq:shock_detection}
\end{equation}

with typical values $C_\text{shock} = 3$ and $C_\text{cd} = 1$.

\subsection{Approximate Riemann Solvers}

For computational efficiency, HLL-type solvers can be used:

\subsubsection{HLLC for SRMHD}

Wave speed estimates:
\begin{align}
S_L &= \min(\lambda_{-,L}^f, \lambda_{-,R}^f) \\
S_R &= \max(\lambda_{+,L}^f, \lambda_{+,R}^f)
\end{align}

Interface pressure (acoustic approximation):
\begin{equation}
p^* = \frac{Z_L p_R + Z_R p_L + Z_L Z_R (v_L - v_R)}{Z_L + Z_R}
\label{eq:hllc_pressure}
\end{equation}

where $Z = \rho_0 h \gamma^2 c_s$ is the relativistic acoustic impedance.

%=============================================================================
\section{Test Problems}
%=============================================================================

\subsection{Relativistic MHD Shock Tubes}

Standard tests include:
\begin{enumerate}
    \item Balsara shock tubes (1-5)
    \item Relativistic rotor problem
    \item Relativistic blast waves with magnetic fields
    \item Relativistic Orszag-Tang vortex
\end{enumerate}

\subsection{Convergence Tests}

Test propagation of:
\begin{itemize}
    \item Relativistic fast magnetosonic waves
    \item Relativistic Alfv\'en waves
    \item Relativistic slow magnetosonic waves
\end{itemize}

Verify second-order convergence ($\epsilon \propto h^2$).

%=============================================================================
\section{Conclusion}
%=============================================================================

This document presents a complete mathematical formulation for Special Relativistic
Magnetohydrodynamics using Godunov SPH, derived from first principles. The key features are:

\begin{enumerate}
    \item \textbf{First-principles derivation}: Starting from special relativity and Maxwell's
        equations, through stress-energy tensor construction
    \item \textbf{Complete primitive recovery}: Full derivation of the nonlinear inversion
        problem including magnetic field effects, with Newton-Raphson algorithm
    \item \textbf{Wave speed derivation}: From linearized equations and dispersion relation
    \item \textbf{Shock conditions}: Complete Rankine-Hugoniot derivation for SRMHD
    \item \textbf{Convolution-based SPH discretization}: Ensuring conservation
    \item \textbf{Decomposition strategy}: Riemann solver for fast shocks, MOC for Alfv\'en waves
    \item \textbf{Powell correction terms}: For divergence-free magnetic field maintenance
\end{enumerate}

This formulation provides a foundation for implementing high-accuracy relativistic MHD
simulations suitable for astrophysical applications including relativistic jets,
magnetized neutron star mergers, and accretion flows around black holes.

%=============================================================================
\appendix
%=============================================================================

\section{Derivation of Relativistic MHD Stress Tensor: Detailed Steps}
\label{app:stress_tensor}

\subsection{Starting Point: Separate Contributions}

The total stress-energy tensor is the sum:
\begin{equation}
T^{\mu\nu} = T^{\mu\nu}_\text{fluid} + T^{\mu\nu}_\text{EM}
\end{equation}

\subsection{Fluid Contribution}

In the fluid rest frame:
\begin{equation}
T^{\mu\nu}_\text{fluid,rest} = \text{diag}(\rho_0 c^2 + \rho_0 \epsilon, p, p, p)
\end{equation}

Boosting to lab frame with 4-velocity $u^\mu$:
\begin{equation}
T^{\mu\nu}_\text{fluid} = (\rho_0 c^2 + \rho_0 \epsilon + p)\frac{u^\mu u^\nu}{c^2} + p \eta^{\mu\nu}
= \rho_0 h \frac{u^\mu u^\nu}{c^2} + p \eta^{\mu\nu}
\end{equation}

\subsection{Electromagnetic Contribution with Ideal MHD}

The EM tensor is:
\begin{equation}
T^{\mu\nu}_\text{EM} = \frac{1}{\mu_0}\left(F^{\mu\alpha}F^\nu_{\ \alpha}
- \frac{1}{4}\eta^{\mu\nu}F^{\alpha\beta}F_{\alpha\beta}\right)
\end{equation}

Using $F^{\mu\nu}$ expressed in terms of $b^\mu$ and $u^\mu$ (ideal MHD):
\begin{equation}
F^{\mu\nu} = u^\mu b^\nu - u^\nu b^\mu
\end{equation}

After lengthy algebra:
\begin{equation}
T^{\mu\nu}_\text{EM} = \frac{b^2}{\mu_0}\frac{u^\mu u^\nu}{c^2} + \frac{b^2}{2\mu_0}\eta^{\mu\nu}
- \frac{b^\mu b^\nu}{\mu_0}
\end{equation}

\subsection{Combined Result}

\begin{equation}
T^{\mu\nu} = \left(\rho_0 h + \frac{b^2}{\mu_0}\right)\frac{u^\mu u^\nu}{c^2}
+ \left(p + \frac{b^2}{2\mu_0}\right)\eta^{\mu\nu}
- \frac{b^\mu b^\nu}{\mu_0}
\end{equation}

\section{Primitive Recovery: Detailed Newton-Raphson Implementation}
\label{app:newton}

\subsection{Function Evaluation}

Given trial $Z$:

\textbf{Step 1:} Compute auxiliary quantities
\begin{align}
\xi &= Z + \frac{B^2}{4\pi} \\
\tilde{\xi} &= \xi + \frac{B^2}{4\pi} = Z + \frac{B^2}{2\pi}
\end{align}

\textbf{Step 2:} Compute $v^2$ from momentum constraint
\begin{equation}
v^2 = \frac{S^2 \tilde{\xi}^2 + (\bm{S} \cdot \bm{B})^2 \xi^2 [(\tilde{\xi}/\xi)^2 - 1]}
{\tilde{\xi}^4}
\end{equation}

After simplification with $\tilde{\xi} = \xi + B^2/(4\pi)$:
\begin{equation}
v^2 = \frac{S^2}{\tilde{\xi}^2} + \frac{(\bm{S} \cdot \bm{B})^2 B^2}{4\pi \xi^2 \tilde{\xi}^2}
\left(1 + \frac{B^2}{4\pi\xi}\right)
\end{equation}

\textbf{Step 3:} Compute Lorentz factor
\begin{equation}
W = \frac{1}{\sqrt{1 - v^2}}
\end{equation}

Check: if $v^2 \geq 1$, state is unphysical.

\textbf{Step 4:} Compute $\bm{v} \cdot \bm{B}$ and $(b^0)^2$
\begin{align}
\bm{v} \cdot \bm{B} &= \frac{(\bm{S} \cdot \bm{B})(4\pi\xi + B^2)}{4\pi\xi\tilde{\xi}} \\
(b^0)^2 &= W^2 (\bm{v} \cdot \bm{B})^2
\end{align}

\textbf{Step 5:} Compute pressure
\begin{equation}
p = Z - Q - \frac{B^2(2W^2-1)}{8\pi W^2} + \frac{(b^0)^2}{8\pi W^2}
\end{equation}

Check: if $p < 0$, apply floor $p = p_\text{floor}$.

\textbf{Step 6:} Compute density
\begin{equation}
\rho_0 = D/W
\end{equation}

\textbf{Step 7:} Compute $Z$ from equation of state
\begin{equation}
Z_\text{eos} = DW + \frac{\Gamma W^2}{\Gamma-1}p
\end{equation}

\textbf{Step 8:} Residual
\begin{equation}
f(Z) = Z - Z_\text{eos}
\end{equation}

\subsection{Derivative Computation}

For Newton-Raphson, we need $f'(Z) = df/dZ$.

Using chain rule through all intermediate variables:
\begin{equation}
f'(Z) = 1 - \frac{dZ_\text{eos}}{dZ} = 1 - \frac{\partial Z_\text{eos}}{\partial W}\frac{dW}{dZ}
- \frac{\partial Z_\text{eos}}{\partial p}\frac{dp}{dZ}
\end{equation}

The individual derivatives are computed by differentiating the expressions in
Step 1-7 with respect to $Z$.

\section{Gauss-Legendre Quadrature for Rarefaction Integrals}
\label{app:quadrature}

For the rarefaction integral:
\begin{equation}
I = \int_{p_a}^{p_b} f(p) dp
\end{equation}

Use Gauss-Legendre quadrature with transformation to $[-1, 1]$:
\begin{equation}
I \approx \frac{p_b - p_a}{2} \sum_{k=1}^{N_G} w_k f\left(\frac{p_a + p_b}{2} + \frac{p_b - p_a}{2} t_k\right)
\end{equation}

Typically $N_G = 16$ points provide sufficient accuracy.

\section{Conservation Properties}
\label{app:conservation}

The antisymmetric form of the SPH equations ensures:

\textbf{Momentum conservation:}
\begin{equation}
\sum_i m_i \frac{d\bm{S}_i}{dt} = 0
\end{equation}

\textbf{Energy conservation:}
\begin{equation}
\sum_i m_i \frac{de_i}{dt} = 0
\end{equation}

These hold exactly when Powell corrections are not applied. With Powell terms,
conservation is sacrificed to maintain numerical stability in low-$\beta$ plasmas.

\section{Limiting Cases and Verification}
\label{app:limits}

\subsection{Non-Relativistic Limit}

As $v/c \to 0$:
\begin{align}
\gamma &\to 1 \\
h &\to c^2 + \epsilon + p/\rho_0 \to c^2 \\
b^0 &\to 0 \\
b^2 &\to B^2
\end{align}

The stress tensor reduces to the classical MHD form.

\subsection{Zero Magnetic Field}

As $\bm{B} \to 0$:
\begin{align}
b^\mu &\to 0 \\
p_\text{tot} &\to p \\
w^* &\to \rho_0 h
\end{align}

The equations reduce to relativistic hydrodynamics.

\subsection{Ultra-Relativistic Limit}

As $\gamma \to \infty$ (and keeping finite proper quantities):
\begin{align}
D &\to \infty \\
\bm{S} &\to \gamma^2 w^* \bm{v} \\
\tau &\to \gamma^2 w^*
\end{align}

The inversion requires careful numerical handling due to large dynamic range.

%=============================================================================
\section{Comparison with Phantom SPH Implementation}
%=============================================================================

This section compares our GSPH formulation with the Phantom SPH code
(Liptai \& Price 2019), which is a widely-used astrophysical SPH code with
general relativistic capabilities.

\subsection{Overview of Phantom's Approach}

Phantom implements GR (general relativistic) hydrodynamics but currently does
\textbf{not} include GRMHD or SRMHD (relativistic MHD). It has:
\begin{itemize}
    \item General relativistic hydrodynamics (curved spacetime)
    \item Non-relativistic MHD (with divergence cleaning)
    \item These are \textit{separate} implementations, not combined
\end{itemize}

\subsection{Stress-Energy Tensor Comparison}

\subsubsection{Phantom (Pure Hydrodynamics)}

Phantom's GR stress-energy tensor (from \texttt{extern\_gr.f90}, line 337):
\begin{equation}
T_{\mu\nu} = \rho h \, u_\mu u_\nu + P g_{\mu\nu}
\label{eq:phantom_tmunu}
\end{equation}

where:
\begin{itemize}
    \item $h = 1 + u + P/\rho$ is the specific enthalpy
    \item $u_\mu$ is the covariant 4-velocity
    \item No magnetic field contribution
\end{itemize}

\subsubsection{Our SRMHD Formulation}

Our complete MHD stress-energy tensor \eqref{eq:Tmunu_MHD}:
\begin{equation}
T^{\mu\nu} = \left(\rho_0 h + \frac{b^2}{\mu_0}\right)\frac{u^\mu u^\nu}{c^2}
+ \left(p + \frac{b^2}{2\mu_0}\right)\eta^{\mu\nu}
- \frac{b^\mu b^\nu}{\mu_0}
\end{equation}

\textbf{Key difference:} We include the magnetic four-vector $b^\mu$ and associated
magnetic pressure/tension terms.

\subsection{Conserved Variables Comparison}

\subsubsection{Phantom's Conserved Variables}

From \texttt{cons2primsolver.f90} (lines 76-127):
\begin{align}
\rho^* &= \sqrt{g} \, \rho \, U^0 \quad \text{(conserved density)} \\
p_i &= U^0 h \, g_{i\mu} v^\mu \quad \text{(covariant momentum)} \\
E &= U^0 h \, g_{\mu\nu} v^\mu v^\nu + \frac{1+u}{U^0} \quad \text{(total energy)}
\end{align}

where $U^0 = \gamma$ is the time component of the 4-velocity.

\subsubsection{Our Conserved Variables}

From \eqref{eq:conserved_variables}:
\begin{align}
D &= \rho_0 \gamma \quad \text{(rest-mass density)} \\
S^i &= w^* \gamma^2 v^i - \frac{b^0 b^i}{4\pi} \quad \text{(momentum with B)} \\
\tau &= w^* \gamma^2 - p_\text{tot} - \frac{(b^0)^2}{4\pi} - D \quad \text{(energy with B)} \\
B^i &\quad \text{(magnetic field)}
\end{align}

\textbf{Key difference:} Our conserved momentum and energy include magnetic field
coupling through $b^0$ and $b^2$ terms.

\subsection{Primitive Variable Recovery Comparison}

\subsubsection{Phantom's Algorithm}

Phantom uses Newton-Raphson iteration on the enthalpy $h$ (lines 181-238 of
\texttt{cons2primsolver.f90}):

\textbf{Root-finding function:}
\begin{equation}
f = 1 + \frac{\Gamma}{\Gamma-1}\frac{P}{\rho} - h_\text{old}
\end{equation}

\textbf{Newton-Raphson derivative:}
\begin{equation}
\frac{df}{dh} = -1 + \frac{\Gamma}{\Gamma-1}\left(1 - \frac{p_\mu p^\mu \cdot P}
{h^3 W^2 \rho}\right)
\end{equation}

\textbf{Lorentz factor from momentum:}
\begin{equation}
W = \sqrt{1 + \frac{p_\mu p^\mu}{h^2}}
\end{equation}

\subsubsection{Our Algorithm (Section 6)}

We iterate on $Z = \rho_0 h W^2$ with the constraint:
\begin{equation}
f(Z) = Z - Z_\text{computed}(Z) = 0
\end{equation}

where $Z_\text{computed}$ is obtained from the velocity constraint and EOS.

\textbf{Key difference:} Our algorithm must handle the magnetic field coupling:
\begin{itemize}
    \item $b^0 = W (\bm{v} \cdot \bm{B})$ couples velocity to B-field
    \item $b^2 = (B^2 + (b^0)^2)/W^2$ appears in energy equation
    \item Momentum includes $-b^0 b^i/(4\pi)$ term
\end{itemize}

This creates a more complex nonlinear system requiring careful handling of the
magnetic contributions at each iteration step.

\subsection{Non-Relativistic MHD in Phantom}

Phantom's non-relativistic MHD (from \texttt{force.F90}, lines 2091-2113) provides
a useful reference for MHD stress tensor implementation:

\textbf{Wave speed:}
\begin{equation}
v_\text{wave} = \sqrt{c_s^2 + v_A^2}, \quad v_A^2 = \frac{B^2}{\rho}
\label{eq:phantom_wave}
\end{equation}

\textbf{MHD stress tensor components:}
\begin{align}
S_{xx} &= S_{xx} - \frac{m_j B_x B_x}{\rho^2} \\
S_{xy} &= S_{xy} - \frac{m_j B_x B_y}{\rho^2} \\
&\vdots
\end{align}

\textbf{Total pressure:}
\begin{equation}
P_\text{tot} = \frac{P + P_\text{rad}}{\rho^2} + \sigma_\text{iso} + \frac{1}{2}\frac{B^2}{\rho^2}
\end{equation}

This corresponds to the non-relativistic limit of our formulation.

\subsection{Key Algorithmic Differences}

\begin{table}[h]
\centering
\begin{tabular}{|l|l|l|}
\hline
\textbf{Aspect} & \textbf{Phantom} & \textbf{Our GSPH} \\
\hline
Physics scope & GR hydro OR Newt. MHD & SRMHD (combined) \\
\hline
Riemann solver & None (standard SPH) & Exact/HLLC for shocks \\
\hline
MHD treatment & Separate module & Integrated in Riemann \\
\hline
Alfv\'en waves & Artificial viscosity & Method of Characteristics \\
\hline
Divergence cleaning & Hyperbolic $\psi$ & Powell source terms \\
\hline
Evolved B-field & $B/\rho$ & $B/D$ (relativistic) \\
\hline
Wave speed & $\sqrt{c_s^2 + v_A^2}$ & Full magnetosonic formula \\
\hline
Con2Prim variable & Enthalpy $h$ & $Z = \rho h W^2$ \\
\hline
\end{tabular}
\caption{Comparison of key algorithmic choices}
\end{table}

\subsection{Advantages of Our GSPH Formulation}

\begin{enumerate}
    \item \textbf{Unified SRMHD:} Magnetic fields are fully integrated into the
        relativistic framework, not treated separately.

    \item \textbf{Riemann solver:} Captures shocks without artificial viscosity,
        providing more accurate shock profiles.

    \item \textbf{MOC for Alfv\'en waves:} The Method of Characteristics properly
        handles the transverse wave dynamics rather than relying on numerical
        dissipation.

    \item \textbf{Full wave speeds:} We use the complete relativistic magnetosonic
        dispersion relation rather than a simplified sum.

    \item \textbf{Stress decomposition:} Following Iwasaki \& Inutsuka (2011),
        we decompose the MHD stress into compressive and Alfv\'enic parts for
        more accurate treatment.
\end{enumerate}

\subsection{Lessons from Phantom}

Several aspects of Phantom's implementation are worth incorporating:

\begin{enumerate}
    \item \textbf{Metric flexibility:} Phantom's \texttt{metric\_tools} module
        provides generic metric handling that could extend our SR formulation to GR.

    \item \textbf{Entropy-based evolution:} Option to evolve entropy instead of
        total energy can improve accuracy in some regimes.

    \item \textbf{Hyperbolic divergence cleaning:} Their $\psi$-based cleaning
        (Tricco \& Price 2012) is an alternative to Powell terms.

    \item \textbf{$B/\rho$ evolution:} Evolving $B/\rho$ (Newtonian) or $B/D$
        (relativistic) naturally preserves the frozen-in condition.
\end{enumerate}

%=============================================================================
\begin{thebibliography}{99}

\bibitem{Kitajima2025} Kitajima, K., Inutsuka, S., \& Seno, I. (2025).
Special Relativistic Smoothed Particle Hydrodynamics Based on Riemann Solver.
\textit{Journal of Computational Physics}, arXiv:2510.18251.

\bibitem{Iwasaki2011} Iwasaki, K., \& Inutsuka, S. (2011).
Smoothed particle magnetohydrodynamics with a Riemann solver and the method of characteristics.
\textit{MNRAS}, 418, 1668.

\bibitem{Inutsuka2002} Inutsuka, S. (2002).
Reformulation of Smoothed Particle Hydrodynamics with Riemann Solver.
\textit{J. Comput. Phys.}, 179, 238.

\bibitem{Pons2000} Pons, J. A., Mart\'i, J. M., \& M\"uller, E. (2000).
The exact solution of the Riemann problem with non-zero tangential velocities in relativistic hydrodynamics.
\textit{J. Fluid Mech.}, 422, 125.

\bibitem{MignioneBodo2005} Mignone, A., \& Bodo, G. (2005).
An HLLC Riemann solver for relativistic flows.
\textit{MNRAS}, 364, 126.

\bibitem{Noble2006} Noble, S. C., Gammie, C. F., McKinney, J. C., \& Del Zanna, L. (2006).
Primitive Variable Solvers for Conservative General Relativistic Magnetohydrodynamics.
\textit{ApJ}, 641, 626.

\bibitem{MignoneMcKinney2007} Mignone, A., \& McKinney, J. C. (2007).
Equation of state in relativistic magnetohydrodynamics: variable versus constant adiabatic index.
\textit{MNRAS}, 378, 1118.

\bibitem{Price2012} Price, D. J. (2012).
Smoothed particle magnetohydrodynamics - IV. Using the vector potential.
\textit{MNRAS}, 423, L45.

\bibitem{Anile1989} Anile, A. M. (1989).
\textit{Relativistic Fluids and Magneto-fluids}.
Cambridge University Press.

\bibitem{Rezzolla2013} Rezzolla, L., \& Zanotti, O. (2013).
\textit{Relativistic Hydrodynamics}.
Oxford University Press.

\bibitem{LiptaiPrice2019} Liptai, D., \& Price, D. J. (2019).
General relativistic smoothed particle hydrodynamics.
\textit{MNRAS}, 485, 819-842.

\bibitem{TriccoPrice2012} Tricco, T. S., \& Price, D. J. (2012).
Constrained hyperbolic divergence cleaning for smoothed particle magnetohydrodynamics.
\textit{J. Comput. Phys.}, 231, 7214-7236.

\bibitem{Dedner2002} Dedner, A., Kemm, F., Kr\"oner, D., Munz, C.-D., Schnitzer, T., \& Wesenberg, M. (2002).
Hyperbolic Divergence Cleaning for the MHD Equations.
\textit{J. Comput. Phys.}, 175, 645-673.

\bibitem{Powell1999} Powell, K. G., Roe, P. L., Linde, T. J., Gombosi, T. I., \& De Zeeuw, D. L. (1999).
A Solution-Adaptive Upwind Scheme for Ideal Magnetohydrodynamics.
\textit{J. Comput. Phys.}, 154, 284-309.

\bibitem{EvansHawley1988} Evans, C. R., \& Hawley, J. F. (1988).
Simulation of magnetohydrodynamic flows: A constrained transport method.
\textit{ApJ}, 332, 659-677.

\bibitem{Gundlach2005} Gundlach, C., Calabrese, G., Hinder, I., \& Mart\'in-Garc\'ia, J. M. (2005).
Constraint damping in the Z4 formulation and harmonic gauge.
\textit{Class. Quantum Grav.}, 22, 3767-3773.

\bibitem{PriceMonaghan2004} Price, D. J., \& Monaghan, J. J. (2004).
Smoothed Particle Magnetohydrodynamics - I. Algorithm and tests in one dimension.
\textit{MNRAS}, 348, 123-138.

\end{thebibliography}

\end{document}
